\documentclass[12pt,a4paper,oneside]{report}

\usepackage[a4paper,top=27mm,bottom=27mm,left=20mm,right=20mm,headheight=15pt,headsep=6mm,footskip=10mm]{geometry}
\usepackage{fontspec}
\setmainfont{Latin Modern Roman}
\setsansfont{Latin Modern Sans}
\setmonofont{Latin Modern Mono}
\usepackage{microtype}
\usepackage{graphicx}
\usepackage{booktabs}
\usepackage{tabularx}
\renewcommand{\tabularxcolumn}[1]{>{\raggedright\arraybackslash}p{#1}}
\usepackage{longtable}
\usepackage{array}
\usepackage{multirow}
\newcolumntype{P}[1]{>{\raggedright\arraybackslash}p{#1}}
\usepackage[table]{xcolor}
\usepackage{tikz}
\usetikzlibrary{arrows.meta,positioning,shapes.geometric,shapes.multipart,fit,calc,backgrounds}
\usepackage{eso-pic}
\usepackage{fancyhdr}
\usepackage{titlesec}
\usepackage{enumitem}
\usepackage{caption}
\usepackage{subcaption}
\usepackage{float}
\usepackage{needspace}
\usepackage{listings}
\usepackage{amsmath}
\usepackage{cite}
\renewcommand{\citepunct}{], [}
\renewcommand{\citedash}{]--[}
\usepackage{hyperref}
\usepackage{url}
\usepackage{setspace}
\usepackage{multicol}
\usepackage[titles]{tocloft}
\setlength{\cftbeforechapskip}{2.5pt}
\setlength{\cftbeforesecskip}{0pt}
\renewcommand{\cftchapfont}{\bfseries\color{teal}}
\renewcommand{\cftchappagefont}{\bfseries}
\setlength{\cftfigindent}{0pt}\setlength{\cfttabindent}{0pt}
\setlength{\cftfignumwidth}{2.4em}\setlength{\cfttabnumwidth}{2.4em}
\setlength{\cftbeforefigskip}{0pt}\setlength{\cftbeforetabskip}{0pt}
\renewcommand{\cftdotsep}{2.5}

\definecolor{ink}{HTML}{000000}
\definecolor{teal}{HTML}{17365D}
\definecolor{rust}{HTML}{365F91}
\definecolor{paper}{HTML}{FFFFFF}
\definecolor{pale}{HTML}{EDF3F8}
\definecolor{line}{HTML}{7E8793}
\definecolor{green}{HTML}{17365D}
\definecolor{linkblue}{HTML}{1F5AA6}
% LostLink interface palette (static/css/tokens.css, light theme)
\definecolor{llpage}{HTML}{F5F3F0}
\definecolor{lltext}{HTML}{211F1E}
\definecolor{llaccent}{HTML}{D6533B}
\definecolor{llaccentink}{HTML}{B83E2B}
\definecolor{llfound}{HTML}{827B75}
\definecolor{llsoft}{HTML}{FFF1ED}

\hypersetup{colorlinks=true,linkcolor=linkblue,urlcolor=linkblue,citecolor=linkblue,pdfauthor={Saanvi Grover, Prankit Goel, Aditya Sharma, Aaradhya Kulashri},pdftitle={LostLink Project Report}}
\setlength{\parindent}{0pt}
\setlength{\parskip}{3pt}
\setstretch{1.00}
\renewcommand{\arraystretch}{1.18}
\setlist[itemize]{leftmargin=5mm,itemsep=1pt,topsep=2pt,parsep=0pt}
\setlist[enumerate]{leftmargin=6mm,itemsep=1pt,topsep=2pt,parsep=0pt}
\captionsetup{font=small,labelfont=bf,labelsep=period,skip=5pt}

% Floats: let text fill the page around figures, and never ship a mostly empty float page.
\renewcommand{\topfraction}{0.9}
\renewcommand{\bottomfraction}{0.8}
\renewcommand{\textfraction}{0.07}
\renewcommand{\floatpagefraction}{0.85}
\setcounter{topnumber}{3}
\setcounter{bottomnumber}{2}
\setcounter{totalnumber}{5}
\setlength{\textfloatsep}{10pt plus 3pt minus 3pt}
\setlength{\floatsep}{8pt plus 2pt minus 2pt}
\setlength{\intextsep}{8pt plus 2pt minus 2pt}

\newif\ifcoverpage
\coverpagetrue
% Page frame on every page (single on the cover, double elsewhere).
\AddToHook{shipout/foreground}{%
  \begin{tikzpicture}[remember picture,overlay]
    \draw[ink,line width=0.55pt]
      ($(current page.north west)+(9mm,-9mm)$) rectangle
      ($(current page.south east)+(-9mm,9mm)$);
    \ifcoverpage\else
      \draw[ink!45,line width=0.25pt]
        ($(current page.north west)+(11mm,-11mm)$) rectangle
        ($(current page.south east)+(-11mm,11mm)$);
    \fi
  \end{tikzpicture}%
}

\pagestyle{fancy}
\fancyhf{}
\fancyhead[L]{\fontsize{9.5}{11}\selectfont\textsc{LostLink Project Report}}
\fancyhead[R]{\fontsize{9.5}{11}\selectfont\nouppercase{\leftmark}}
\fancyfoot[C]{\small\thepage}
\renewcommand{\headrulewidth}{0.3pt}
\renewcommand{\footrulewidth}{0pt}
\fancypagestyle{plain}{\fancyhf{}\fancyfoot[C]{\small\thepage}\renewcommand{\headrulewidth}{0pt}}

% Every chapter starts on a new page.
\setlength{\emergencystretch}{2em}

\titleformat{\chapter}[display]
  {\normalfont\color{ink}}
  {\fontsize{10}{12}\selectfont\sffamily\bfseries\color{teal}\MakeUppercase{Chapter \thechapter}}
  {2pt}
  {\fontsize{21}{25}\selectfont\bfseries}
  [\vspace{1pt}{\color{teal}\rule{20mm}{0.6pt}}]
\titlespacing*{\chapter}{0pt}{0pt}{10pt}
\titleformat{\section}{\fontsize{14}{17}\selectfont\bfseries\color{teal}}{\thesection}{0.55em}{}
\titlespacing*{\section}{0pt}{8pt}{3pt}
\titleformat{\subsection}{\fontsize{12.5}{15}\selectfont\bfseries\color{ink}}{\thesubsection}{0.55em}{}
\titlespacing*{\subsection}{0pt}{6pt}{2pt}
\setcounter{tocdepth}{1}

\newcommand{\status}[1]{\colorbox{pale}{\textcolor{teal}{\sffamily\bfseries #1}}}
\newcommand{\code}[1]{\texttt{#1}}
\newcommand{\assetroot}{assets}
\newcommand{\shotdir}{\assetroot/shots}
\newcommand{\shotframe}[2][\linewidth]{\fbox{\includegraphics[width=#1]{\shotdir/#2}}}
\newcommand{\swatch}[1]{\tikz[baseline=-0.6ex]\draw[ink!40,fill=#1] (0,-0.12) rectangle (0.55,0.22);}
\setlength{\fboxsep}{0pt}
\setlength{\fboxrule}{0.3pt}

% Numbered IEEE-style entries without the report class's duplicate bibliography heading.
\makeatletter
\renewenvironment{thebibliography}[1]
  {\list{\@biblabel{\@arabic\c@enumiv}}%
    {\settowidth\labelwidth{\@biblabel{#1}}%
     \leftmargin\labelwidth
     \advance\leftmargin\labelsep
     \usecounter{enumiv}%
     \let\p@enumiv\@empty
     \renewcommand\theenumiv{\@arabic\c@enumiv}}%
   \sloppy\clubpenalty4000\widowpenalty4000\sfcode`\.\@m}
  {\def\@noitemerr{\@latex@warning{Empty `thebibliography' environment}}\endlist}
\makeatother

\lstdefinestyle{lostlink}{
  language=Python,
  basicstyle=\ttfamily\footnotesize,
  keywordstyle=\color{ink}\bfseries,
  commentstyle=\color{ink!70},
  stringstyle=\color{ink},
  frame=single,
  rulecolor=\color{line},
  backgroundcolor=\color{paper},
  breaklines=true,
  showstringspaces=false,
  columns=fullflexible,
  xleftmargin=3mm,
  framexleftmargin=2mm,
  aboveskip=6pt,
  belowskip=6pt
}

% Shared TikZ styles
\tikzset{
  box/.style={draw=ink,rounded corners=1.5mm,fill=paper,align=center,font=\small,inner sep=2.2mm},
  pbox/.style={box,fill=pale},
  dec/.style={diamond,draw=rust,fill=rust!8,aspect=2.4,align=center,font=\scriptsize,inner sep=1.2mm},
  term/.style={draw=ink,rounded corners=4mm,fill=pale,align=center,font=\small\bfseries,inner sep=2mm},
  arr/.style={-{Latex[length=2mm]},thick,color=teal},
  lab/.style={font=\scriptsize,text=ink!80,fill=paper,inner sep=1pt},
}
\begin{document}

% ---------------------------------------------------------------------------
% Title page - follows the supplied Full Stack Development Lab synopsis cover.
% ---------------------------------------------------------------------------
\begin{titlepage}
\thispagestyle{empty}
\begingroup
\begin{center}
\vspace*{2mm}
{\fontsize{18}{22}\selectfont\bfseries PROJECT REPORT\par}
\vspace{3mm}
{\fontsize{12}{15}\selectfont on\par}
\vspace{3mm}
{\fontsize{18}{23}\selectfont\bfseries LOSTLINK\par}
\vspace{2mm}
{\fontsize{15}{19}\selectfont\bfseries AN AI-ENABLED FULL-STACK PLATFORM FOR\\INTELLIGENT LOST-AND-FOUND ITEM MATCHING\\AND RECOVERY\par}
\vspace{6mm plus 1fill}
\includegraphics[height=45mm]{\assetroot/synopsis-media/image1.jpg}\par
\vspace{6mm plus 1fill}
{\fontsize{12.5}{16}\selectfont\bfseries Full Stack Development Lab (24B15CS314)\par}
\vspace{1mm}
{\fontsize{11.5}{15}\selectfont Project Based Learning (PBL)\par}
\vspace{1mm}
{\fontsize{11.5}{15}\selectfont Odd Semester 2026, Academic Session 2026--27\par}
\end{center}

\vspace{6mm plus 1fill}
\noindent\hspace*{10mm}%
{\fontsize{11}{17}\selectfont
\begin{tabular}[t]{@{}l@{}}
{\fontsize{12}{15}\selectfont\bfseries SUBMITTED TO}\\[2mm]
Dr. Sonal\\
Dr. Shobhit Tyagi
\end{tabular}}\hfill
{\fontsize{11}{17}\selectfont
\begin{tabular}[t]{@{}l@{\hspace{8mm}}r@{}}
\multicolumn{2}{@{}l}{\fontsize{12}{15}\selectfont\bfseries SUBMITTED BY}\\[2mm]
Saanvi Grover & 2408030010\\
Prankit Goel & 2408030011\\
Aditya Sharma & 2408030019\\
Aaradhya Kulashri & 2408030026
\end{tabular}}\hspace*{10mm}

\vspace{6mm plus 1fill}
\begin{center}
{\fontsize{12}{15}\selectfont\bfseries
DEPARTMENT OF COMPUTER SCIENCE AND ENGINEERING\\
JAYPEE INSTITUTE OF INFORMATION TECHNOLOGY,\\
NOIDA (U.P.) NOVEMBER 2026\par}
\end{center}
\vspace{2mm}
\endgroup
\end{titlepage}
\coverpagefalse

\pagenumbering{roman}
\setcounter{page}{1}

% ---------------------------------------------------------------------------
\chapter*{Abstract}
\addcontentsline{toc}{chapter}{Abstract}
\markboth{Abstract}{Abstract}
LostLink is a full-stack campus lost-and-found platform that replaces scattered messages and notice-board posts with a searchable, moderated and auditable recovery workflow. Students report lost or found objects with structured details and an optional photograph, and a moderator reviews every report before it is published. A hybrid matching engine compares approved lost and found reports on six signals and stores a per-signal explanation of every score. A similarity score proposes a lead; it never proves ownership. Recovery runs through a separate human-in-the-loop claim process: private verification questions and evidence, an independent moderator decision recorded on an evidence checklist, controlled disclosure of contact details, and a physical handover confirmed with a one-time code held by the owner. Claim approved is not item returned.

\textbf{Key contributions}
\begin{itemize}
  \item \textbf{Stable, explainable matching:} engine \code{synopsis-v3} scores title and description with a fixed, versioned lexical model, so a pair's score never depends on which other reports exist; ResNet-50 image features, zone, category, graded date-and-time chronology and colour sets complete six signals, each shown with its contribution and the share of evidence actually compared.
  \item \textbf{Calibrated image similarity:} image cosine is rescaled above each model's measured chance level, so unrelated photographs earn no points.
  \item \textbf{Ownership integrity:} an approved claim reserves the item (one per item, enforced by the database); disputes never release it, rejected claimants can request reconsideration, and the handover completes only when the owner's one-time code is entered.
  \item \textbf{Campus identity and privacy:} verified campus email accounts, password recovery, rate limiting, independent moderation, sanitised uploads, blurred sensitive photographs and an append-only audit trail.
  \item \textbf{Operations:} three roles, analytics, CSV/PDF exports, first-run initialisation without demo data, and production configuration checks.
  \item \textbf{Polished interface:} Black + Vermilion design system with Light/Dark/System themes, disc-overlap match visualisation and a phone tab bar.
\end{itemize}

\textbf{Technology:} Django 6.1 on Python 3.13, server-rendered templates with progressive JavaScript, SQLite for development and PostgreSQL in production.

\textbf{Verification:} on the merged repository (commit \code{3eda809}), all 364 automated tests pass on SQLite, with one optional pretrained-model test skipped by design. An 11-step first-run test, a 318-check browser sweep, a 65-check mobile layout test and a complete end-to-end recovery journey ending in a one-time-code handover also pass.

\textbf{Limitation:} matching accuracy has not yet been measured on independently photographed, human-reviewed campus reports; the bundled pilot set supports engineering checks only.

\textbf{Keywords:} lost and found, Django, explainable matching, evidence coverage, image-similarity calibration, role-based access control, one-time-code handover, ResNet-50, human verification.

% ---------------------------------------------------------------------------
\clearpage
\chapter*{Table of Contents}
\markboth{Contents}{Contents}
\begingroup
\fontsize{9.6}{11}\selectfont
\setlength{\columnsep}{7mm}
\begin{multicols}{2}
\makeatletter\@starttoc{toc}\makeatother
\end{multicols}
\endgroup

% ---------------------------------------------------------------------------
\clearpage
\chapter*{List of Figures and Tables}
\addcontentsline{toc}{chapter}{List of Figures and Tables}
\markboth{List of Figures and Tables}{List of Figures and Tables}
\begingroup
\fontsize{7.9}{8.9}\selectfont
\makeatletter
\setlength{\columnsep}{6mm}
\noindent\begin{minipage}[t]{0.505\textwidth}
\textbf{\color{teal}Figures}\par\vspace{1pt}
\@starttoc{lof}
\end{minipage}\hfill
\begin{minipage}[t]{0.465\textwidth}
\textbf{\color{teal}Tables}\par\vspace{1pt}
\@starttoc{lot}
\end{minipage}
\makeatother
\endgroup
\vspace{3pt}
{\color{teal}\bfseries Abbreviations}\par
{\footnotesize\raggedright
\textbf{AJAX} Asynchronous JavaScript and XML; \textbf{API} Application Programming Interface; \textbf{CLIP} Contrastive Language--Image Pre-training; \textbf{CNN} Convolutional Neural Network; \textbf{CRUD} Create, Read, Update, Delete; \textbf{CSRF} Cross-Site Request Forgery; \textbf{DFD} Data Flow Diagram; \textbf{ERD} Entity Relationship Diagram; \textbf{HSTS} HTTP Strict Transport Security; \textbf{MRR} Mean Reciprocal Rank; \textbf{MVT} Model--View--Template; \textbf{OTP} One-Time Password (handover code); \textbf{ORM} Object--Relational Mapping; \textbf{RBAC} Role-Based Access Control; \textbf{SMTP} Simple Mail Transfer Protocol; \textbf{TF--IDF} Term Frequency--Inverse Document Frequency; \textbf{UI} User Interface.\par}

\clearpage
\pagenumbering{arabic}
\setcounter{page}{1}

% ===========================================================================
\chapter{Introduction}

\section{Project Overview}
LostLink gives the campus a single recovery desk.
\begin{itemize}
  \item A person files a \textbf{lost} or \textbf{found} report with a title, description, category, campus zone, date and an optional photograph.
  \item A moderator reviews the report; only approved reports become searchable and eligible for matching.
  \item The matching engine compares every eligible lost--found pair and ranks plausible connections, showing how much each signal contributed.
  \item A person who believes an item is theirs opens a private \textbf{claim}; an independent moderator checks the evidence, and approval \textbf{reserves} the item and reveals contact details.
  \item At the meeting the owner creates a \textbf{one-time code} and shares it only after receiving the item; the person handing it over enters the code, which closes both reports and every dependent claim and match.
\end{itemize}

\section{Problem Statement}
Campus recovery today depends on fragmented channels, and each has a gap:
\begin{itemize}
  \item \textbf{Scattered:} class groups, informal messages, security desks and paper notices never meet in one place.
  \item \textbf{Unsearchable:} free-text posts are quickly buried and cannot be filtered by place, date or category.
  \item \textbf{Inconsistent:} two people describe the same object in different words, and photographs are often missing.
  \item \textbf{Unsafe:} posting phone numbers publicly exposes personal information, and anyone can claim an item.
  \item \textbf{Unaccountable:} nobody records who handed what to whom, so disputes cannot be resolved.
\end{itemize}
A useful system must improve discovery without treating algorithmic similarity as proof of ownership.

\begin{table}[H]
\centering\small
\caption[Existing channels compared with LostLink]{Existing recovery channels compared with LostLink}
\begin{tabularx}{\textwidth}{X cccc}
\toprule
Capability & Chat groups & Notice board & Security desk & LostLink\\
\midrule
Single searchable record & -- & -- & Partial & Yes\\
Filter by zone, date and category & -- & -- & -- & Yes\\
Automatic ranking of likely pairs & -- & -- & -- & Yes\\
Contact details kept private & -- & -- & Yes & Yes\\
Ownership verified before handover & -- & -- & Partial & Yes\\
Audit trail of decisions & -- & -- & Partial & Yes\\
Available from a phone at any time & Yes & -- & -- & Yes\\
\bottomrule
\end{tabularx}
\end{table}

\section{Objectives}
\begin{enumerate}
  \item Centralise lost and found reports in a structured, searchable database.
  \item Rank likely pairs using explainable text, visual, location, date and metadata signals.
  \item Require moderation before a report becomes public or eligible for matching.
  \item Verify ownership with private questions, private evidence and an independent staff decision, and confirm the physical return with the owner's one-time code.
  \item Limit accounts to verified campus members and protect contact details, private answers, verification secrets and sensitive photographs.
  \item Give administrators auditable analytics, exports and decision history.
\end{enumerate}

\section{Scope of the Project}
\textbf{In scope:} verified campus accounts, password recovery and three roles (Student, Moderator, Administrator); reports with one photograph, moderation, search and hybrid matching; claims with private questions and evidence, disputes, reconsideration and the one-time-code handover; in-app notifications, analytics, CSV and PDF exports.

\textbf{Out of scope:} native mobile apps, payments, GPS tracking, SMS, email beyond account verification and password reset, and any automatic ownership approval.

\section{Motivation}
The project answers a practical campus need and an engineering question: can multimodal similarity narrow a noisy search while keeping people in control of the consequential decision? LostLink separates \emph{discovery} from \emph{verification}: the \textbf{machine proposes} a ranked, explained lead; a \textbf{moderator verifies} ownership from private evidence; the \textbf{owner's one-time code confirms} the physical handover.

\section{Organisation of the Report}
Chapters 2--3 cover requirements and the technology stack; Chapter 4 the architecture, data model and UML diagrams; Chapter 5 the matching engine, ownership workflow, one-time-code handover and security; Chapters 6--10 the modules, interface, testing, results and conclusions.

% ===========================================================================
\chapter{Requirement Analysis}
\textbf{In this chapter:}
\begin{itemize}[beginpenalty=10000]
  \item who uses LostLink and what each stakeholder needs;
  \item functional and non-functional requirements, with their status;
  \item use cases, platform requirements, role capabilities and acceptance evidence.
\end{itemize}

\section{Stakeholders}
\begin{itemize}
  \item Five groups have a stake in campus recovery; two of them are the students on either side of a lost item.
  \item Each need below maps to a feature described later in this report.
\end{itemize}
\begin{table}[H]
\centering\small
\caption{Stakeholders and their needs}
\begin{tabularx}{\textwidth}{P{30mm}X X}
\toprule
Stakeholder & Need & How LostLink serves it\\
\midrule
Student who lost an item & Find it quickly without trawling every channel & One search wall, ranked matches and notifications\\
Student who found an item & Hand it back safely to the real owner & Private claims, moderator-verified ownership and a code-confirmed handover\\
Moderator (staff) & Review reports and decide claims fairly & Review queue, evidence view and decision log\\
Administrator & Run the service and understand its use & User and taxonomy management, analytics, exports\\
Institution & Only campus members; protect student data; resolve disputes & Verified campus email, RBAC, blurred sensitive photos, append-only audit trail\\
\bottomrule
\end{tabularx}
\end{table}

\section{Functional Requirements}
\begin{itemize}
  \item Twelve functional requirements cover the whole journey, from verified registration to closing a case.
  \item All twelve are implemented and verified by automated tests (Chapter~8).
\end{itemize}
\begin{table}[H]
\centering\small
\caption{Functional requirements and implemented response}
\begin{tabularx}{\textwidth}{>{\bfseries}P{13mm}X P{27mm}}
\toprule
ID & Requirement & Status\\
\midrule
FR-01 & Register a student account, verify the campus email with an expiring single-use link, sign in, reset a forgotten password. & \status{Implemented}\\
FR-02 & Create, edit, view, withdraw and flag lost/found reports. & \status{Implemented}\\
FR-03 & Review reports independently (never one's own) and publish only approved, active records. & \status{Implemented}\\
FR-04 & Search by keyword, type, category, zone, status, date and sort order. & \status{Implemented}\\
FR-05 & Compute stable, explainable lost--found similarity with evidence coverage and heuristic tiers. & \status{Implemented}\\
FR-06 & Start, question, answer, decide (with an evidence checklist), dispute, reconsider and withdraw a claim. & \status{Implemented}\\
FR-07 & Notify relevant users without duplicating unread events. & \status{Implemented}\\
FR-08 & Administer users, taxonomy, analytics and filtered exports. & \status{Implemented}\\
FR-09 & Close every dependent claim, match and report when a case ends. & \status{Implemented}\\
FR-10 & Confirm the physical handover with a one-time code held by the owner. & \status{Implemented}\\
FR-11 & Keep verification secrets out of public text; blur sensitive photographs. & \status{Implemented}\\
FR-12 & Initialise a fresh production database without demo accounts. & \status{Implemented}\\
\bottomrule
\end{tabularx}
\end{table}

\section{Use Case Specifications}
\begin{itemize}
  \item Eight use cases cover the whole recovery journey, from filing a report to running the service.
  \item Each one names its actor and the outcome stored in the database.
\end{itemize}
\begin{table}[H]
\centering\small
\caption{Principal use cases}
\begin{tabularx}{\textwidth}{>{\bfseries}P{10mm}P{29mm}P{25mm}X}
\toprule
ID & Use case & Actor & Main flow and outcome\\
\midrule
UC-01 & File a report & Student & Fill the form; report saved as \emph{pending} and queued for review\\
UC-02 & Review a report & Moderator & Approve (publish and match) or reject with a reason; reporter notified\\
UC-03 & Inspect a match & Student & Open a report; see candidates with per-signal scores\\
UC-04 & Claim an item & Student & Submit a statement; finder notified; questions follow\\
UC-05 & Decide a claim & Moderator & Complete the evidence checklist; approve (item reserved, contacts shared) or reject with a reason\\
UC-06 & Hand over & Owner + other party & Owner creates a one-time code after receiving the item; the other enters it; case closed\\
UC-07 & Raise a dispute & Participant & Recorded beside the claim; reservation kept; handover paused for an independent moderator\\
UC-08 & Manage the service & Administrator & Roles, taxonomy, analytics and CSV/PDF exports\\
\bottomrule
\end{tabularx}
\end{table}

\section{Non-Functional Requirements}
\begin{itemize}
  \item Quality requirements describe how well the system behaves, not what it does.
  \item Security, privacy and integrity carry the most weight, because the system handles personal belongings and contact details.
\end{itemize}
\begin{table}[H]
\centering\small
\caption{Quality requirements}
\begin{tabularx}{\textwidth}{P{27mm}X}
\toprule
Quality & Requirement\\
\midrule
Security & Verified campus accounts, rate-limited sign-in, CSRF protection, server-side roles, object checks, deny-by-default media and safe exports.\\
Privacy & Answers, expected answers, private details, evidence and contacts are never public; sensitive photographs are blurred for others.\\
Integrity & Independent review; one reserved claim per item; disputes never release it; a hashed one-time code confirms the return.\\
Usability & Responsive interface, keyboard access, no-JavaScript fallbacks, reduced-motion support and clear empty and error states.\\
Reliability & Database constraints, row locks, idempotent re-matching, recorded feature-extraction status and tolerant input parsing.\\
Maintainability & Modular Django apps, migrations, shared templates, design tokens, tests and documented configuration.\\
Performance & Cached image and text features, pagination, indexed relations and production PostgreSQL support.\\
Explainability & Every match stores its component scores, evidence status and coverage, chronology, model manifest and engine version.\\
\bottomrule
\end{tabularx}
\end{table}

\section{Hardware and Software Requirements}
\begin{itemize}
  \item LostLink runs on an ordinary laptop for development; the image model is the only heavy component.
\end{itemize}
\begin{table}[H]
\centering\small
\caption{Hardware and software requirements}
\begin{tabularx}{\textwidth}{P{34mm}X}
\toprule
Item & Requirement\\
\midrule
Processor and memory & Dual-core CPU and 8 GB RAM (lexical text + histogram profile); 16 GB recommended for ResNet-50; GPU optional\\
Storage & About 2 GB free, plus space for uploaded photographs and model weights\\
Runtime & Python 3.11 or later (developed on 3.13) and Django 6.x \cite{python,django}\\
Libraries & Pillow, NumPy, scikit-learn \cite{pillow,sklearn,sklearnpaper}; ReportLab for PDF \cite{reportlab}\\
Optional ML profile & PyTorch, TorchVision, Sentence Transformers and CLIP \cite{pytorch,pytorchpaper}\\
Database & PostgreSQL 17 in production; SQLite for local development \cite{postgres}\\
Email & An SMTP account in production for verification and password-reset links (console output in development)\\
Client & Any standards-compliant desktop or mobile browser\\
\bottomrule
\end{tabularx}
\end{table}

\section{User Requirements}
\begin{itemize}
  \item Capabilities grow with the role; administrators can do everything moderators can.
  \item Acceptance is judged by the visible result, the permission boundary and the stored state together.
\end{itemize}
\begin{table}[H]
\centering\small
\caption{Role-to-capability matrix}
\begin{tabularx}{\textwidth}{P{34mm}XXX}
\toprule
Capability & Student & Moderator & Administrator\\
\midrule
Search and browse & Yes (also signed out) & Yes & Yes\\
Report and claim & Yes, once email is verified & Yes & Yes\\
Hold or enter a handover code & Owner / person handing over & Only as a participant & Only as a participant\\
Review reports and claims & No & Yes, except own cases & Yes, except own cases\\
Resolve flags, disputes, exceptions & No & Yes, if independent & Yes, if independent\\
Manage users and taxonomy & No & No & Yes\\
Analytics and exports & No & Limited review context & Yes\\
\bottomrule
\end{tabularx}
\end{table}

\begin{table}[H]
\centering\small
\caption{User acceptance evidence}
\begin{tabularx}{\textwidth}{P{28mm}X P{33mm}}
\toprule
Area & Acceptance criterion & Verification evidence\\
\midrule
Student access & A student can manage only their own reports and their own claims. & Ownership and role tests\\
Moderation & Unapproved reports stay non-public and are never matched. & Workflow integration tests\\
Matching & Every proposed pair shows its component scores and missing evidence. & Scoring and persistence tests\\
Recovery & Contact details appear only after an independent, checklist-backed approval. & End-to-end claim test\\
Handover & Only the owner's valid, unexpired, single-use code marks the item returned. & 15 handover tests + browser journey\\
Case closure & A confirmed return closes the claimant's own report, its matches and its claims. & Regression tests\\
Identity & Unverified accounts cannot report or claim; tokens expire and work once. & Verification tests\\
Administration & Exports omit private answers and neutralise formula-like CSV values. & Export security tests\\
\bottomrule
\end{tabularx}
\end{table}

% ===========================================================================
\chapter{Technology Stack}
\textbf{In this chapter:}
\begin{itemize}[beginpenalty=10000]
  \item the six-layer stack from browser to server;
  \item frontend, backend and database choices and why they were made;
  \item the tools and custom management commands used to build and run the project.
\end{itemize}

\section{Overview}
\begin{itemize}
  \item Six layers, each with a single responsibility, from the browser down to the server.
  \item No frontend build step and no external API dependency in the core workflow.
  \item Machine-learning models are optional and swappable through \code{matching/profile.json}.
\end{itemize}
\begin{figure}[!htbp]
\centering
\begin{tikzpicture}[layer/.style={draw=ink,rounded corners=1.2mm,minimum width=158mm,minimum height=8.5mm,align=left,font=\footnotesize,anchor=north west,inner xsep=3mm,text width=152mm},
  tag/.style={font=\small\bfseries,text=teal,anchor=west}]
\node[layer,fill=pale] (l1) at (0,0) {\hspace{36mm}Django templates; 5 CSS files with design tokens; 3 JS files};
\node[layer] (l2) at (0,-1.05) {\hspace{36mm}Django 6.1: routing, forms, sessions, CSRF, role and verification decorators, rate limits};
\node[layer,fill=pale] (l3) at (0,-2.1) {\hspace{36mm}accounts, items, matching, claims, notifications, dashboard};
\node[layer] (l4) at (0,-3.15) {\hspace{36mm}Versioned lexical text model, ResNet-50 (PyTorch), optional CLIP and Sentence-BERT};
\node[layer,fill=pale] (l5) at (0,-4.2) {\hspace{36mm}Django ORM; SQLite (dev); PostgreSQL 17 (prod); protected media};
\node[layer] (l6) at (0,-5.25) {\hspace{36mm}Gunicorn; WhiteNoise; SMTP for account email; Render web service with a persistent disk};
\foreach \n/\t in {l1/Presentation,l2/Web framework,l3/Domain apps,l4/Machine learning,l5/Data,l6/Serving}
  \node[tag] at ($(\n.west)+(1.5mm,0)$) {\t};
\end{tikzpicture}
\caption{LostLink technology stack, from browser to server}
\end{figure}

\section{Frontend Technologies}
\begin{itemize}
  \item \textbf{Templates:} semantic HTML rendered by Django, with a shared shell and reusable partials.
  \item \textbf{Styles:} five modular CSS files (tokens, base, components, layout, landing); custom properties define the palette, both themes, type, spacing and motion.
  \item \textbf{Scripts:} three JavaScript files add debounced AJAX search, notification polling, live previews, confirmation dialogs, match-overlap animation and theme synchronisation \cite{mdn}.
  \item \textbf{Progressive enhancement:} every core workflow still works with JavaScript disabled.
\end{itemize}

\section{Backend and Database Technologies}
\begin{itemize}
  \item \textbf{Django} supplies routing, forms, authentication, the ORM, templates, CSRF protection and test utilities \cite{django}.
  \item \textbf{Six apps} hold the business rules: \code{accounts}, \code{items}, \code{matching}, \code{claims}, \code{notifications} and \code{dashboard}.
  \item \textbf{Gunicorn and WhiteNoise} serve the application and its static files in production \cite{whitenoise}.
  \item \textbf{SQLite} is the zero-configuration development database; \textbf{PostgreSQL} is the deployment target with stronger concurrency \cite{postgres}.
  \item \textbf{Migrations, foreign keys, indexes and partial unique constraints} keep both schemas aligned and enforce recovery invariants: one reserved claim per item, one open dispute per claim and one active handover code per claim.
  \item \textbf{Django's cache} backs the rate limits for sign-in, registration, verification, password reset and handover codes; \textbf{SMTP} delivers verification and reset links.
\end{itemize}

\begin{table}[H]
\centering\small
\caption{Technology stack summary}
\begin{tabularx}{\textwidth}{P{28mm}P{48mm}X}
\toprule
Layer & Technology & Purpose\\
\midrule
Presentation & HTML, CSS, JavaScript & Responsive Black + Vermilion interface and AJAX search\\
Application & Django 6.1, Python 3.13 & Request handling, authentication, forms and ORM\\
Data & SQLite / PostgreSQL 17 & Development and production persistence\\
Text similarity & Versioned lexical model (term weighting after TF--IDF \cite{tfidf}) & Title + description vectors over a fixed vocabulary; cosine similarity\\
Image similarity & ResNet-50 \cite{resnet}; CLIP \cite{clip} optional & Versioned, calibrated visual feature vectors\\
Media & Pillow & Decoding, re-encoding without metadata, blurred previews\\
Security & \code{secrets}, HMAC-SHA-256, Django signing & One-time handover codes, verification and reset tokens\\
Exports & ReportLab, CSV & Administrative summaries and privacy-safe records\\
Serving & Gunicorn, WhiteNoise & Production WSGI server and static files\\
Testing & Django TestCase, Playwright & Unit and integration tests, browser workflows\\
\bottomrule
\end{tabularx}
\end{table}

\section{Development Tools and Commands}
\begin{itemize}
  \item \textbf{Git and GitHub} for version control and pull-request review \cite{lostlinkrepo}.
  \item \textbf{Docker Compose} provisions a local PostgreSQL server for compatibility testing.
  \item \textbf{No frontend build step:} the project runs directly from \code{manage.py}.
\end{itemize}

\begin{table}[H]
\centering\small
\caption{Custom management commands}
\begin{tabularx}{\textwidth}{P{53mm}X}
\toprule
Command & Purpose\\
\midrule
\code{seed\_taxonomy} & Create the default categories and campus zones on a fresh database (idempotent)\\
\code{seed\_demo} & Demo users, reports and matches (refuses to run in production)\\
\code{mark\_accounts\_verified} & Mark restored or legacy accounts as email-verified\\
\code{sanitize\_photos} & Re-encode stored photos without metadata; flag likely sensitive ones\\
\code{purge\_claim\_evidence} & Delete evidence files past the retention period\\
\code{build\_text\_model} & Compile the versioned lexical text release (immutable once built)\\
\code{rebuild\_matches} & Rescore every match after an engine, model or taxonomy change\\
\code{retry\_features}, \code{matching\_status} & Retry failed feature extraction; report feature and match health\\
\code{build\_pilot}, \code{evaluate\_hybrid} & Build the pilot set; compare model combinations using production scoring\\
\code{evaluate\_matches}, \code{retune} & Evaluate a labelled split; tune on validation data only\\
\code{validate\_research\_dataset} & Check provenance, splits and leakage before research use\\
\code{monitor\_matching} & Store a score-distribution snapshot and flag drift\\
\bottomrule
\end{tabularx}
\end{table}

% ===========================================================================
\chapter{System Design}
\textbf{In this chapter:}
\begin{itemize}[beginpenalty=10000]
  \item architecture, request lifecycle and deployment design;
  \item use case, data flow, entity relationship and activity diagrams;
  \item state machines, sequence diagrams, the database schema and the URL map.
\end{itemize}

\section{System Architecture}
\begin{itemize}
  \item \textbf{Layered and server-rendered:} requests pass through Django middleware and role checks before views call forms, domain services and ORM queries.
  \item \textbf{Separated matching:} feature extraction, calibration and scoring live in their own service layer and can run without HTTP.
  \item \textbf{Central business rules:} every report status change goes through one lifecycle function and every claim action through one service module, so views, admin pages and commands cannot disagree.
  \item \textbf{Deny-by-default media:} photographs are re-encoded on upload and served through a permission-aware view; anything without an explicit rule is refused.
\end{itemize}

\begin{figure}[!htbp]
\centering
\begin{tikzpicture}[scale=0.66,every node/.append style={transform shape},node distance=6mm]
\node[pbox,minimum width=150mm,minimum height=9mm] (browser) {Responsive browser UI\\Django templates + CSS tokens + progressive JavaScript};
\node[box,below=of browser,minimum width=150mm,minimum height=9mm] (web) {Django request layer: URLs, middleware, sessions, CSRF, decorators, forms and views};
\node[box,below=8mm of web,minimum width=46mm,minimum height=11mm] (match) {Matching services\\features, calibration, scoring};
\node[box,left=8mm of match,minimum width=46mm,minimum height=11mm] (domain) {Domain rules\\accounts, item lifecycle, claim services};
\node[box,right=8mm of match,minimum width=46mm,minimum height=11mm] (admin) {Operations\\notifications, moderation, analytics};
\node[box,below=12mm of match,minimum width=60mm,minimum height=11mm] (db) {Django ORM\\SQLite / PostgreSQL};
\node[box,right=10mm of db,minimum width=50mm,minimum height=11mm] (media) {Protected media storage\\sanitised photos, private evidence};
\draw[arr] (browser) -- (web);
\draw[arr] (web.south -| domain) -- (domain);
\draw[arr] (web) -- (match);
\draw[arr] (web.south -| admin) -- (admin);
\draw[arr] (domain.south) -- (db.north west);
\draw[arr] (match) -- (db);
\draw[arr] (admin.south) -- (db.north east);
\draw[arr] (match.south east) -- (media.north west);
\end{tikzpicture}
\caption{Layered architecture of LostLink}
\end{figure}

\section{Request Lifecycle}
Every page follows Django's Model--View--Template flow, with security checks before any data is read.
\begin{figure}[!htbp]
\centering
\begin{tikzpicture}[scale=0.92,every node/.append style={transform shape},node distance=5mm and 6mm]
\node[term] (req) {HTTP request};
\node[box,right=of req] (mw) {Middleware\\session, CSRF,\\security headers};
\node[box,right=of mw] (url) {URL\\resolver};
\node[dec,right=of url] (auth) {Verified, role and\\ownership OK?};
\node[box,right=of auth] (view) {View\\form + service};
\node[box,below=12mm of view] (orm) {ORM\\query / save};
\node[box,left=of orm] (tpl) {Template\\render};
\node[term,left=of tpl] (resp) {HTML or JSON\\response};
\node[box,above=7mm of auth,fill=rust!8] (deny) {403 / 429 page, sign-in or verify-email redirect (keeps \code{?next=})};
\draw[arr] (req) -- (mw); \draw[arr] (mw) -- (url); \draw[arr] (url) -- (auth);
\draw[arr] (auth) -- node[lab,above]{yes} (view);
\draw[arr] (auth) -- node[lab,right]{no} (deny);
\draw[arr] (view) -- (orm); \draw[arr] (orm) -- (tpl); \draw[arr] (tpl) -- (resp);
\end{tikzpicture}
\caption{Request lifecycle through Django's Model--View--Template layers}
\end{figure}

\section{Deployment Architecture}
\begin{itemize}
  \item Production runs as one Render web service connected to a managed PostgreSQL database.
  \item Photographs and model weights live on a persistent disk, because the service file system is reset on every deploy.
  \item An SMTP account sends verification and password-reset email; the app refuses to start with an unsafe secret key, host list or CSRF origin.
\end{itemize}
\begin{figure}[!htbp]
\centering
\begin{tikzpicture}[scale=0.88,every node/.append style={transform shape},node distance=6mm and 10mm]
\node[pbox,minimum width=30mm,minimum height=12mm] (client) {Browser\\(phone or laptop)};
\node[draw=teal,dashed,rounded corners=2mm,minimum width=78mm,minimum height=46mm,right=14mm of client,yshift=-8mm] (svc) {};
\node[font=\scriptsize\bfseries,text=teal,anchor=north west] at (svc.north west) {\,Render web service};
\node[box,minimum width=62mm] (gun) at ($(svc.center)+(0,10mm)$) {Gunicorn (WSGI) $\rightarrow$ Django app};
\node[box,minimum width=29mm,below left=6mm and -31mm of gun] (wn) {WhiteNoise\\static files};
\node[box,minimum width=29mm,below right=6mm and -31mm of gun] (ml) {Matching engine\\ResNet-50, lexical text};
\node[box,minimum width=36mm,right=16mm of svc.north east,anchor=north west] (pg) {PostgreSQL\\(managed database)};
\node[box,minimum width=36mm,below=10mm of pg] (disk) {Persistent disk\\\code{/var/data/media}\\model cache};
\node[box,minimum width=36mm,below=8mm of disk] (smtp) {SMTP server\\verification, reset};
\draw[arr,<->] (client) -- node[lab,above]{HTTPS} (svc.west |- client);
\draw[arr] (gun.east) -- node[lab,above,pos=0.3]{\code{DATABASE\_URL}} (pg.west);
\draw[arr] (ml.east) -- (disk.west);
\draw[arr] (svc.south east) -- (smtp.west);
\draw[arr] (gun) -- (wn); \draw[arr] (gun) -- (ml);
\end{tikzpicture}
\caption{Production deployment design for Render (a free demonstration instance runs separately)}
\end{figure}
\begin{itemize}
  \item \textbf{Build:} \code{build.sh} installs requirements and runs \code{collectstatic}; before each deploy \code{migrate} and \code{seed\_taxonomy --if-empty} run, so a fresh database is usable without demo data.
  \item \textbf{Configuration:} \code{DJANGO\_DEBUG=false}, a private secret key of at least 50 characters, \code{DATABASE\_URL}, exact allowed hosts, https CSRF origins, \code{DJANGO\_EMAIL\_*} and optionally \code{LOSTLINK\_ALLOWED\_EMAIL\_DOMAINS}.
  \item \textbf{Storage:} photographs need a persistent disk, which is why the free tier (no disk) is unsuitable.
  \item \textbf{After an engine upgrade:} run \code{rebuild\_matches} once.
\end{itemize}

\section{Use Case Diagram}
\begin{itemize}
  \item \textbf{Three actors:} Student, Moderator and Administrator; the Administrator inherits every Moderator use case.
  \item \textbf{Students} report, search, claim, and hand items over with the owner's one-time code.
  \item \textbf{Moderators} review reports and decide claims (never their own), and resolve disputes and documented handover exceptions.
  \item \textbf{Administrators} also manage users, taxonomy, analytics and exports.
\end{itemize}
\begin{figure}[!htbp]
\centering
\begin{tikzpicture}[scale=0.78,every node/.append style={transform shape},
  actor/.style={draw=ink,circle,minimum size=9mm,font=\small\bfseries,fill=paper},
  role/.style={font=\small\bfseries,anchor=east},
  use/.style={draw=ink,ellipse,minimum width=44mm,minimum height=9mm,inner sep=0.5mm,align=center,font=\small,fill=pale},
  link/.style={draw=ink!70},
  inherit/.style={-{Triangle[open,length=3mm,width=3mm]},draw=ink,thick}
]
\draw[ink,rounded corners=2mm] (2.4,5.0) rectangle (15.4,-4.6);
\node[font=\scriptsize\bfseries,anchor=north west] at (2.6,4.85) {LostLink System};
\node[actor] (student) at (0,3.0) {S};
\node[role] at (-0.7,3.0) {Student};
\node[actor] (mod) at (0,0.0) {M};
\node[role] at (-0.7,0.0) {Moderator};
\node[actor] (admin) at (0,-3.0) {A};
\node[role] at (-0.7,-3.0) {Administrator};
\node[use] (report) at (6.0,3.9) {Report or edit item};
\node[use] (search) at (12.0,3.9) {Search and inspect};
\node[use] (claim) at (6.0,2.3) {Claim, answer, appeal};
\node[use] (notify) at (12.0,2.3) {Hand over with one-time code};
\node[use] (review) at (6.0,0.7) {Review reports};
\node[use] (verify) at (12.0,0.7) {Decide claims};
\node[use] (handover) at (9.0,-0.9) {Resolve disputes, exceptions};
\node[use] (users) at (6.0,-3.0) {Manage users and taxonomy};
\node[use] (analytics) at (12.0,-3.0) {Analytics and exports};
\draw[link] (student.east) -- (9.0,3.0);
\draw[link] (2.9,3.0) |- (report.west);
\draw[link] (2.9,3.0) |- (claim.west);
\draw[link] (9.0,3.0) |- (search.west);
\draw[link] (9.0,3.0) |- (notify.west);
\draw[link] (mod.east) -- (9.0,0.0);
\draw[link] (2.9,0.0) |- (review.west);
\draw[link] (9.0,0.0) |- (verify.west);
\draw[link] (9.0,0.0) -- (handover.north);
\draw[link] (admin.south) -- (0,-4.1) -- (12.0,-4.1);
\draw[link] (6.0,-4.1) -- (users.south);
\draw[link] (12.0,-4.1) -- (analytics.south);
\draw[inherit] (admin.north) -- node[right,font=\scriptsize]{inherits role} (mod.south);
\end{tikzpicture}
\caption{Principal use cases by role}
\end{figure}

\section{Data Flow Diagram}
\begin{itemize}
  \item \textbf{Four processes:} report and search, match and explain, verify and recover, moderate and analyse.
  \item \textbf{One data store} (D1) holds reports, matches, claims, notifications and decisions.
  \item \textbf{Two external entities:} campus users and staff; only approved reports flow into matching.
\end{itemize}
\begin{figure}[!htbp]
\centering
\begin{tikzpicture}[
  ext/.style={draw=ink,rectangle,minimum width=27mm,minimum height=11mm,fill=paper,align=center,font=\small},
  proc/.style={draw=teal,rectangle,rounded corners=2mm,minimum width=28mm,minimum height=12mm,fill=pale,align=center,font=\small},
  store/.style={draw=rust,double,rectangle,minimum width=42mm,minimum height=10mm,fill=paper,align=center,font=\small},
  flow/.style={-{Latex[length=2mm]},thick,color=ink!72},
]
\node[ext] (user) at (-7.4,2.4) {Campus user};
\node[proc] (reports) at (-2.9,2.4) {1.0 Report\\and search};
\node[proc] (matching) at (1.6,2.4) {2.0 Match\\and explain};
\node[proc] (claims) at (6.1,2.4) {3.0 Verify\\and recover};
\node[store] (data) at (1.6,-0.2) {D1 Application data};
\node[proc] (ops) at (-2.9,-2.8) {4.0 Moderate\\and analyse};
\node[ext] (staff) at (6.1,-2.8) {Moderator / Admin};
\draw[flow] (user) -- node[lab,above=1pt]{query} (reports);
\draw[flow] (reports) -- node[lab,above=1pt]{approved} (matching);
\draw[flow] (matching) -- node[lab,above=1pt]{ranked lead} (claims);
\draw[flow] (claims) -- node[lab,right]{claim status} (staff);
\draw[flow] (staff) -- node[lab,above]{decision} (ops);
\draw[flow] (ops) -- node[lab,left]{review state} (reports);
\draw[flow] (reports.south) |- node[lab,pos=0.75,above]{write} (data.west);
\draw[flow] (matching.south) -- node[lab,right]{read / write} (data.north);
\draw[flow] (claims.south) |- node[lab,pos=0.75,above]{case data} (data.east);
\end{tikzpicture}
\caption{Level-0 data flow diagram}
\end{figure}

\section{Entity Relationship Diagram}
\begin{itemize}
  \item \textbf{Item} is the hub: it belongs to a User, a Category and a Location.
  \item \textbf{Match} links one lost and one found Item; retired rows stay as history. \textbf{ItemFeatures} records how each report's text and photo were processed.
  \item \textbf{Claim} links a claimant to an Item and owns its questions, statement revisions, private evidence, disputes and handover codes.
  \item \textbf{ReportFlag} and \textbf{ModerationDecision} provide the moderation and append-only audit trail.
\end{itemize}
\begin{figure}[!htbp]
\centering
\begin{tikzpicture}[scale=0.74,every node/.append style={transform shape},
  entity/.style={draw=ink,rounded corners=1mm,fill=paper,minimum width=28mm,minimum height=11mm,align=center,font=\small\bfseries},
  wideentity/.style={entity,minimum width=38mm},
  rel/.style={draw=teal,thick},
  rlabel/.style={fill=paper,inner sep=1.2pt,font=\scriptsize,align=center,text=ink!78}
]
\node[font=\small\bfseries,text=teal] at (-2.0,6.2) {Reporting and matching};
\node[font=\small\bfseries,text=teal] at (9.0,6.2) {Recovery and audit};
\node[entity] (user) at (-6.5,2) {User};
\node[entity] (item) at (-2,2) {Item};
\node[entity] (cat) at (2.4,4.2) {Category};
\node[entity] (loc) at (2.4,0.4) {Location};
\node[entity] (match) at (-2,-1.6) {Match};
\node[entity] (flag) at (-6.5,-1.6) {ReportFlag};
\node[entity] (feat) at (-2,4.3) {ItemFeatures};
\node[entity] (claim) at (7,2) {Claim};
\node[wideentity] (question) at (12.2,5.0) {VerificationQuestion};
\node[wideentity] (revision) at (12.2,3.3) {ClaimRevision};
\node[wideentity] (evidence) at (12.2,1.6) {ClaimEvidence};
\node[wideentity] (dispute) at (12.2,-0.1) {ClaimDispute};
\node[wideentity] (handover) at (12.2,-1.8) {HandoverSession};
\node[wideentity] (decision) at (7,-3.9) {ModerationDecision};
\node[entity] (note) at (2.4,-2.2) {Notification};
\draw[rel] (user) -- node[rlabel,above=1pt]{1 : N} (item);
\draw[rel] (item.north east) -- node[rlabel,above left]{N : 1} (cat.west);
\draw[rel] (item.south east) -- node[rlabel,below left]{N : 1} (loc.west);
\draw[rel] (item) -- node[rlabel,right]{lost / found pair} (match);
\draw[rel] (item.south west) -- node[rlabel,above left]{1 : N} (flag.north east);
\draw[rel] (item) -- node[rlabel,right]{1 : 1} (feat);
\draw[rel] (item.east) ++(0,0.25) -- node[rlabel,above,pos=0.62]{1 : N claims} ($(claim.west)+(0,0.25)$);
\draw[rel] (user.north) -- ++(0,3.0) -| node[rlabel,pos=0.4,above]{1 : N claimants} (claim.north);
\draw[rel] (claim.east) -- node[rlabel,pos=0.55]{1 : N} (question.west);
\draw[rel] (claim.east) -- (revision.west);
\draw[rel] (claim.east) -- (evidence.west);
\draw[rel] (claim.east) -- node[rlabel,pos=0.55]{1 : N} (dispute.west);
\draw[rel] (claim.east) -- node[rlabel,pos=0.6]{1 : N} (handover.west);
\draw[rel] (claim) -- node[rlabel,right]{audit 1 : N} (decision);
\draw[rel] (claim.south west) -- node[rlabel,above left]{events} (note.north east);
\end{tikzpicture}
\caption{Entity relationship model (simplified)}
\end{figure}

\section{Activity Flow}
\begin{itemize}
  \item Two decision points: moderator approval and whether a plausible pair exists.
  \item Rejected reports are returned with a reason; unmatched reports stay searchable.
  \item The flow ends only when the owner's one-time code confirms the handover and closes every dependent record.
\end{itemize}
\begin{figure}[!htbp]
\centering
\begin{tikzpicture}[scale=0.8,every node/.append style={transform shape},node distance=5mm,
  flowstep/.style={box,minimum width=52mm,minimum height=8mm}]
\node[term] (start) {Start};
\node[flowstep,below=of start] (file) {Student files report};
\node[flowstep,below=of file] (validate) {Form and image validation};
\node[dec,below=of validate] (review) {Approved?};
\node[flowstep,below left=8mm and 10mm of review] (reject) {Reject with reason;\\reporter notified};
\node[flowstep,below right=8mm and 10mm of review] (publish) {Publish and run matching};
\node[dec,below=of publish] (candidate) {Plausible pair?};
\node[flowstep,below left=8mm and 2mm of candidate] (search) {Remain searchable};
\node[flowstep,below right=8mm and -14mm of candidate] (case) {Owner opens private claim};
\node[flowstep,below=of case] (questions) {Questions, evidence, checklist decision};
\node[flowstep,below=of questions] (return) {Item reserved; owner's code entered};
\node[flowstep,below=of return] (close) {Close both reports, matches and claims};
\node[term,below=of close] (stop) {End};
\draw[arr] (start)--(file); \draw[arr] (file)--(validate); \draw[arr] (validate)--(review);
\draw[arr] (review)-|node[lab,above,pos=0.25]{No}(reject); \draw[arr] (review)-|node[lab,above,pos=0.25]{Yes}(publish);
\draw[arr] (publish)--(candidate); \draw[arr] (candidate)-|node[lab,above,pos=0.25]{No}(search); \draw[arr] (candidate)--node[lab,right]{Yes}(case);
\draw[arr] (case)--(questions); \draw[arr] (questions)--(return); \draw[arr] (return)--(close); \draw[arr] (close)--(stop);
\end{tikzpicture}
\caption{Report-to-recovery activity flow}
\end{figure}

\section{State Machines}
Reports and claims move through explicit states. Every report transition goes through \code{items.lifecycle}, which rejects any move not drawn below; claim services lock the item and then the claim before changing either, so a repeated or concurrent request cannot skip a step.
\begin{figure}[!htbp]
\centering
\begin{tikzpicture}[
  st/.style={draw=ink,rounded corners=2mm,fill=pale,minimum width=24mm,minimum height=8mm,font=\small\bfseries},
  end/.style={st,fill=paper,double}]
\node[st] (p) at (0,0) {Pending};
\node[st] (a) at (4.8,0) {Approved};
\node[st] (c) at (9.6,0) {Claimed};
\node[end] (r) at (14.4,0) {Returned};
\node[st] (x) at (1.6,-2.8) {Rejected};
\node[end] (w) at (6.4,-2.8) {Withdrawn};
\node[end] (rm) at (11.2,-2.8) {Removed};
\draw[arr] (p) -- node[lab,above,align=center]{moderator\\approves} (a);
\draw[arr] (a) -- node[lab,above,align=center]{claim\\opened} (c);
\draw[arr] (c) -- node[lab,above,align=center]{owner's code\\verified} (r);
\draw[arr] (c.north) to[bend right=30] node[lab,above]{claims closed or withdrawn} (a.north);
\draw[arr] (p.south) ++(-0.3,0) -- node[lab,left]{rejected} (x.north);
\draw[arr] (x.east) to[bend right=25] node[lab,right,pos=0.7]{edited} (p.south east);
\draw[arr] (a.south) -- node[lab,left]{withdrawn} (w.north);
\draw[arr] (c.south) -- node[lab,right]{flag upheld} (rm.north);
\end{tikzpicture}
\caption{Report (item) lifecycle}
\end{figure}

\begin{figure}[!htbp]
\centering
\begin{tikzpicture}[
  st/.style={draw=ink,rounded corners=2mm,fill=pale,minimum width=21mm,minimum height=8mm,font=\small\bfseries},
  end/.style={st,fill=paper,double}]
\node[st] (s) at (0,0) {Submitted};
\node[st] (q) at (3.7,0) {Questions};
\node[st] (n) at (7.4,0) {Answered};
\node[st,fill=llsoft] (ap) at (11.1,0) {Approved};
\node[end] (rt) at (14.8,0) {Returned};
\node[end] (rj) at (7.4,-2.7) {Rejected};
\node[end] (wd) at (3.7,-2.7) {Withdrawn};
\node[draw=rust,dashed,rounded corners=2mm,fill=paper,minimum width=36mm,minimum height=8mm,font=\scriptsize,align=center] (d) at (12.9,-2.7) {Open dispute (separate record)\\keeps status; pauses handover};
\draw[arr] (s) -- node[lab,above]{asked} (q);
\draw[arr] (q) -- node[lab,below]{answered} (n);
\draw[arr] (n) -- node[lab,above,align=center]{checklist\\approval} (ap);
\draw[arr] (ap) -- node[lab,above,align=center]{owner's\\code} (rt);
\draw[arr] (n) -- node[lab,left]{rejected} (rj);
\draw[arr] (ap.south west) -- node[lab,below right,pos=0.35]{revoked} (rj.north east);
\draw[arr] (rj.north west) to[bend left=25] node[lab,left,pos=0.4]{reconsider} (n.south west);
\draw[arr] (q.south) -- node[lab,left]{withdrawn} (wd.north);
\draw[arr] (n.north) to[bend right=35] node[lab,above]{follow-up question} (q.north);
\draw[ink!50,dashed] (ap.south east) -- (d.north);
\end{tikzpicture}
\caption[Claim lifecycle]{Claim lifecycle. Approved is the ownership reservation (one per item); approving one claim rejects rival open claims; a dispute is recorded alongside and never releases it.}
\end{figure}

\section{Sequence Diagrams}
\begin{itemize}
  \item \textbf{Figure~\ref{fig:seqreport}} follows a report from submission to the first match notification.
  \item \textbf{Figure~\ref{fig:seqclaim}} follows a claim on a found report from the first statement to the code-confirmed handover; the claimant is the owner, so the claimant holds the code.
  \item Solid arrows are requests or notifications; dashed arrows are responses or disclosures.
\end{itemize}
\begin{figure}[!htbp]
\centering
\begin{tikzpicture}[x=1cm,y=-0.62cm,
  hd/.style={draw=ink,fill=pale,rounded corners=1mm,minimum width=25mm,minimum height=7mm,font=\small\bfseries},
  msg/.style={-{Latex[length=2mm]},thick,color=teal},
  ret/.style={-{Latex[length=2mm]},dashed,color=ink!60},
  ml/.style={font=\scriptsize,fill=paper,inner sep=1pt,above=0pt}]
\foreach \x/\n in {0/Student A,3.6/Django app,7.2/Moderator,10.8/Matching engine,14.4/Student B}{
  \node[hd] at (\x,0) {\n}; \draw[ink!40,dashed] (\x,0.6) -- (\x,13.4);}
\draw[msg] (0,1.5) -- node[ml]{submit lost report} (3.6,1.5);
\draw[ret] (3.6,2.5) -- node[ml]{saved as pending} (0,2.5);
\draw[msg] (3.6,3.5) -- node[ml]{appears in review queue} (7.2,3.5);
\draw[msg] (7.2,4.7) -- node[ml]{approve} (3.6,4.7);
\draw[msg] (3.6,5.9) -- node[ml]{run\_matching(item)} (10.8,5.9);
\draw[ret] (10.8,7.1) -- node[ml]{scored pairs + components} (3.6,7.1);
\draw[msg] (3.6,7.9) -- ++(0.7,0) -- ++(0,0.8) -- (3.6,8.7); \node[font=\scriptsize,anchor=west] at (4.4,8.3) {store Match};
\draw[msg] (3.6,9.5) -- node[ml]{``Report approved'' + match summary} (0,9.5);
\draw[msg] (3.6,10.7) -- node[ml,pos=0.62]{``Similar report'' (strong/moderate only)} (14.4,10.7);
\draw[msg] (14.4,12.0) -- node[ml,pos=0.38]{open report: see score breakdown} (3.6,12.0);
\end{tikzpicture}
\caption{Sequence: report submission, moderation, matching and notification}\label{fig:seqreport}
\end{figure}

\begin{figure}[!htbp]
\centering
\begin{tikzpicture}[x=0.93cm,y=-0.6cm,
  hd/.style={draw=ink,fill=pale,rounded corners=1mm,minimum width=25mm,minimum height=7mm,font=\small\bfseries},
  msg/.style={-{Latex[length=2mm]},thick,color=teal},
  ret/.style={-{Latex[length=2mm]},dashed,color=ink!60},
  ml/.style={font=\scriptsize,fill=paper,inner sep=1pt,above=0pt}]
\foreach \x/\n in {0/Owner (claimant),4.8/LostLink,9.6/Moderator,14.4/Finder}{
  \node[hd] at (\x,0) {\n}; \draw[ink!40,dashed] (\x,0.6) -- (\x,19.4);}
\draw[msg] (0,1.5) -- node[ml]{start claim + statement} (4.8,1.5);
\draw[msg] (4.8,2.6) -- node[ml]{``Someone claimed your item''} (14.4,2.6);
\draw[msg] (14.4,3.8) -- node[ml]{private question (+ expected answer)} (4.8,3.8);
\draw[msg] (4.8,5.0) -- node[ml]{question waiting} (0,5.0);
\draw[msg] (0,6.2) -- node[ml]{answer, evidence photo} (4.8,6.2);
\draw[msg] (4.8,7.4) -- node[ml]{ready for review} (9.6,7.4);
\draw[msg] (9.6,8.6) -- node[ml]{checklist approval (independent)} (4.8,8.6);
\draw[ret] (4.8,9.8) -- node[ml]{item reserved; contacts shared} (0,9.8);
\draw[ret] (4.8,9.8) -- node[ml]{contacts shared} (14.4,9.8);
\node[draw=ink!50,fill=llsoft,rounded corners=1mm,font=\scriptsize,minimum width=140mm] at (7.2,11.1) {they meet; the finder hands the item over and the owner checks it};
\draw[msg] (0,12.4) -- node[ml]{create handover code} (4.8,12.4);
\draw[ret] (4.8,13.6) -- node[ml]{6-digit code, shown once} (0,13.6);
\draw[ret] (0,14.8) -- node[ml,pos=0.3]{tells the code, only after receiving the item} (14.4,14.8);
\draw[msg] (14.4,16.0) -- node[ml]{enter code} (4.8,16.0);
\draw[msg] (4.8,17.0) -- ++(0.7,0) -- ++(0,0.8) -- (4.8,17.8); \node[font=\scriptsize,anchor=west] at (5.6,17.4) {verify HMAC, expiry, attempts};
\draw[ret] (4.8,18.8) -- node[ml]{returned; both reports closed} (0,18.8);
\draw[ret] (4.8,18.8) -- node[ml]{handover completed} (14.4,18.8);
\end{tikzpicture}
\caption{Sequence: ownership verification and handover}\label{fig:seqclaim}
\end{figure}

\section{Database Schema}
\begin{itemize}
  \item Fifteen models across six apps; status fields use enumerated choices.
  \item Partial unique constraints enforce one reserved claim per item, one open dispute per claim and one active handover code per claim.
\end{itemize}
\begin{table}[H]
\centering\scriptsize
\caption{Core database schema}
\begin{tabularx}{\textwidth}{P{28mm}P{50mm}X}
\toprule
Model & Representative fields & Key relationships and constraints\\
\midrule
User & username, email, role, enrolment\_no, phone, email\_verified\_at & Role is student, moderator or admin; verification method recorded (email, reset, admin, legacy)\\
Category / Location & name, icon, related or adjacent entries & Controlled vocabulary; changing related or adjacent entries rescores affected matches\\
Item & type, title, description, image, status, event date and time, brand, colour, private\_details, photo\_sensitive & Reporter, Category and Location keys; status includes withdrawn and removed\\
ItemFeatures & text and image state, model used, attempts, error & One per report; distinguishes a missing photo from a failed extraction\\
Match & lost\_item, found\_item, score, tier, six components, coverage, chronology, status & One row per pair; retired rows keep history; engine version and manifest stored\\
Claim & item, claimant, statement, status, decision note, appeal\_count, handover\_method, returned\_at & At most one approved or returned claim per item (partial unique constraint)\\
VerificationQuestion & text, expected\_answer, answer, asked\_by & Expected answer visible only to the asker and staff\\
ClaimRevision & number, kind, statement, responds\_to\_decision & Immutable statement history: original and each reconsideration\\
ClaimEvidence & file, note, uploaded\_by, purged\_at & Re-encoded private photos; claimant and staff only; purged after retention\\
ClaimDispute & raised\_by, reason, status, outcome, resolution & Never changes the claim status; one open dispute per claim\\
HandoverSession & recipient, provider, salt, otp\_hash, expires\_at, failed\_attempts, status & Code stored as a salted HMAC; one active session per claim\\
Notification & user, item, kind, title, url, read state & Recipient-private; duplicate unread events are reused\\
ReportFlag & item, reporter, reason, resolved\_at & Enters the moderation queue; resolved when the report is withdrawn or decided\\
ModerationDecision & actor, action, item, claim, from\_state, to\_state, note, metadata & Append-only: updates and deletes are refused\\
ScoreSnapshot & histogram, mean, tier counts, manifest & Score-distribution monitoring between engine versions\\
\bottomrule
\end{tabularx}
\end{table}

\section{URL Map}
\begin{itemize}
  \item Every staff and administrator route is guarded by a role decorator on the server.
  \item Search and browse are public; everything that reveals people or private evidence requires sign-in.
\end{itemize}
\begin{table}[H]
\centering\small
\caption{Principal routes and who can open them}
\begin{tabularx}{\textwidth}{P{64mm}X P{25mm}}
\toprule
Route & Page or action & Access\\
\midrule
\code{/} & Landing page (redirects signed-in users home) & Public\\
\code{/accounts/login/}, \code{/register/} & Sign-in and registration (rate limited) & Public\\
\code{/accounts/verify/...}, \code{/password/reset/...} & Email verification and password recovery & Public / signed in\\
\code{/dashboard/} & Student home: matches, claims, alerts & Signed in\\
\code{/report/}, \code{/item/<id>/edit/} & File or edit a report & Verified owner\\
\code{/search/}, \code{/lost/}, \code{/found/} & Search and browse approved reports (JSON at \code{/search/api/}) & Public\\
\code{/item/<id>/} & Report detail with match explanations & Signed in\\
\code{/claims/start/<id>/}, \code{/claims/<id>/} & Start and follow a claim & Verified participants, staff\\
\code{/claims/<id>/reconsider/}, \code{evidence/}, \code{withdraw/} & Appeal, private evidence, withdraw & Claimant\\
\code{/claims/<id>/decide/<decision>/} & Approve (checklist) or reject a claim & Independent staff\\
\code{/claims/<id>/dispute/}, \code{dispute/resolve/} & Raise or resolve a dispute & Parties / independent staff\\
\code{/claims/<id>/handover/code/}, \code{verify/} & Create or enter the one-time code & Owner / person handing over\\
\code{/notifications/} & Notification inbox & Signed in\\
\code{/manage/moderation/} & Review queue and decision log & Staff\\
\code{/manage/analytics/}, \code{/export/csv|pdf/} & Analytics and exports & Administrator\\
\code{/manage/users/}, \code{/taxonomy/} & People, roles and vocabulary & Administrator\\
\code{/media/<path>} & Photographs; blurred preview when sensitive & Authorised users\\
\bottomrule
\end{tabularx}
\end{table}
% ===========================================================================
\chapter{Implementation}
\textbf{In this chapter:}
\begin{itemize}[beginpenalty=10000]
  \item how the frontend, backend, authentication and validation are built;
  \item the matching engine: stable text model, image calibration, colour and chronology rules, evidence coverage;
  \item the ownership workflow and the one-time-code handover;
  \item notifications, security controls and the workflow rules fixed in two logic audits.
\end{itemize}

\section{Frontend Implementation}
\begin{itemize}
  \item \textbf{Shared shell and partials} for branding, icons, report cards, overlap diagrams, match rows, status markers and pagination.
  \item \textbf{Layered CSS:} design tokens $\rightarrow$ base rules $\rightarrow$ components $\rightarrow$ page layouts; the landing page adds its own stylesheet and script.
  \item \textbf{Live search:} filters are serialised into shareable URLs, JSON fragments are fetched with skeleton states, and the plain form still works without JavaScript.
  \item \textbf{Theme controller:} stores Light, Dark or System in local storage and applies it before first paint, so pages never flash the wrong theme.
  \item \textbf{Cache busting:} the \code{\{\% asset \%\}} template tag appends a file version during development.
\end{itemize}

\section{Backend Implementation}
\begin{itemize}
  \item Each bounded context has its own URL module and views.
  \item Views combine role and verification decorators, object-level checks, forms and service functions.
  \item Business rules live in two places only: \code{items/lifecycle.py} (every report status change) and \code{claims/services.py} (every claim action); views lock rows, call them and render.
  \item Matching is split into feature extraction, calibration, scoring and persistence, so it can be tested without HTTP.
  \item Management commands expose maintenance and evaluation without coupling them to web requests.
  \item Status fields use enumerated choices; JSON fields store score breakdowns and model manifests.
  \item PostgreSQL is selected through \code{DATABASE\_URL}; otherwise a local SQLite file is used.
\end{itemize}

\section{Authentication and Authorization}
\begin{itemize}
  \item \textbf{Sessions} backed by a custom user model; self-registration always creates a Student and never signs the user in.
  \item \textbf{Campus identity:} a signed, expiring, single-use link verifies the email address; an optional domain list (\code{LOSTLINK\_ALLOWED\_EMAIL\_DOMAINS}) restricts sign-up to the institution. Existing accounts were marked verified as \emph{legacy} by a data migration.
  \item \textbf{Gates:} reporting, flagging and every claim action need \code{@verified\_required}; staff pages need \code{@moderator\_required}; administration needs \code{@admin\_required}.
  \item \textbf{Recovery and abuse protection:} password reset by expiring single-use link (other sessions signed out); sign-in, registration, resend and reset are rate limited and answer identically whether or not an address exists; sign-out is a POST.
  \item \textbf{Object checks:} only the owner edits a report; only participants and staff open a claim; a role alone never reveals another person's answers or contacts.
  \item \textbf{Return path:} a signed-out visitor sent to sign-in keeps the original address in \code{?next=}.
\end{itemize}

\begin{table}[H]
\centering\small
\caption{Claim actions by participant}
\footnotesize
\begin{tabularx}{\textwidth}{X P{16mm}P{17mm}P{17mm}P{19mm}}
\toprule
Action on a claim & Claimant & Finder (reporter) & Moderator & Admin\\
\midrule
View statement and progress & Yes & Yes & Yes & Yes\\
Ask a question (with private expected answer) & -- & Yes & Yes$^{*}$ & Yes$^{*}$\\
Answer; add private evidence; withdraw & Yes & -- & -- & --\\
Request reconsideration after rejection & Yes & -- & -- & --\\
Approve (checklist) or reject & -- & -- & Yes$^{*}$ & Yes$^{*}$\\
See contact details & After approval & After approval & Yes & Yes\\
Raise a dispute / resolve it & Raise & Raise & Resolve$^{*}$ & Resolve$^{*}$\\
Create / enter the handover code$^{\dagger}$ & Create & Enter & -- & --\\
Record a handover exception & -- & -- & Yes$^{*}$ & Yes$^{*}$\\
\bottomrule
\multicolumn{5}{p{\dimexpr\textwidth-2\tabcolsep}}{\footnotesize $^{*}$ Only if independent of the case. $^{\dagger}$ For a found report; on a lost report the roles swap (the owner always holds the code).}
\end{tabularx}
\end{table}

\section{CRUD Operations and Validation}
\begin{itemize}
  \item Students create and update their own reports; withdrawal is a lifecycle transition, not a deletion, so the audit trail survives.
  \item Administrators create, edit and delete categories and zones; staff approve or reject reports and claims.
  \item Every state-changing request is a CSRF-protected POST; redirects are limited to the same host.
  \item Server-side validation is authoritative even when the browser gives instant feedback.
\end{itemize}

\begin{table}[H]
\centering\small
\caption{Form validation rules}
\begin{tabularx}{\textwidth}{P{30mm}X}
\toprule
Field & Rule\\
\midrule
Title & At least 3 characters\\
Description & At least 10 characters\\
Event date & Cannot be in the future\\
Title, description & Refused if they contain an ID, serial or phone number, a labelled serial/IMEI/roll number or an email address; these belong in the private details field\\
Photograph & JPG, PNG, WEBP or GIF, checked by signature and full decode; under 5 MB and 40 megapixels; HTML/SVG polyglots refused; re-encoded as JPEG without metadata under a random name\\
Evidence photo & Same checks; stored privately, re-encoded, at most five per claim\\
Email & Lower-cased, unique, within the allowed campus domains; verified before reporting or claiming\\
Password & Django validators: length, similarity, common-password and numeric checks\\
Role & Never accepted from the registration form; only administrators change it\\
\bottomrule
\end{tabularx}
\end{table}

\section{Matching Engine}
\subsection{Scoring Model}
\begin{itemize}
  \item \textbf{Trigger:} runs when a report is approved, edited or rebuilt.
  \item \textbf{Candidates:} active, approved reports of the opposite type by another person, whose item is not already reserved by an approved claim (\code{Item.objects.matchable()}).
  \item \textbf{Six signals} $s_i \in [0,1]$ with synopsis weights $w_i$; a signal without evidence on both reports is dropped and the weights renormalised, and the share of weight actually compared is recorded as the \textbf{evidence coverage} $C$:
\end{itemize}
\begin{equation}
S \;=\; \frac{\sum_{i \in A} w_i\, s_i}{\sum_{i \in A} w_i}, \qquad C \;=\; \sum_{i \in A} w_i, \qquad A = \{\, i : \text{evidence for } i \text{ exists on both reports} \,\}.
\end{equation}

\begin{table}[H]
\centering\small
\caption{Signal definitions and weights}
\begin{tabularx}{\textwidth}{P{30mm}c X}
\toprule
Signal & Weight & How $s_i$ is computed\\
\midrule
Title + description & 0.30 & Cosine of fixed, versioned lexical vectors (Section~\ref{sec:text})\\
Photograph & 0.30 & Calibrated cosine of ResNet-50 feature vectors (Section~\ref{sec:calib})\\
Campus zone & 0.15 & 1 same zone; 0.5 adjacent zone; 0 otherwise\\
Category & 0.10 & 1 same category; 0.5 related category; 0 otherwise\\
Date and time & 0.10 & $0.5^{|\Delta d|/10}$ proximity, adjusted by the chronology grade (Section~\ref{sec:chrono})\\
Brand / colour & 0.05 & Brand aliases; colour families compared as sets (Section~\ref{sec:attr})\\
\bottomrule
\end{tabularx}
\end{table}

\begin{itemize}
  \item \textbf{Storage:} pairs scoring at least 0.40 are stored with all six components, each component's evidence status, the coverage, the chronology grade and the engine version (\code{synopsis-v3}).
  \item \textbf{Heuristic tiers:} strong similarity above 0.85, moderate 0.70--0.85, weak below 0.70. These are documented thresholds, not calibrated probabilities.
  \item \textbf{Caps:} coverage below 0.70 allows at most moderate and below 0.50 only weak; suspicious chronology caps at moderate, contradictory at weak.
  \item \textbf{Notifications:} only strong and moderate leads notify, once per pair; reruns never notify again.
\end{itemize}

\begin{figure}[!htbp]
\centering
\begin{tikzpicture}[scale=0.84,every node/.append style={transform shape},node distance=5mm and 8mm,
  st/.style={box,minimum width=46mm,minimum height=8mm}]
\node[term] (s) {Report approved};
\node[st,below=of s] (cand) {Load opposite-type, approved,\\unreserved reports};
\node[st,below=of cand] (feat) {Reuse cached lexical and image vectors\\(record feature status)};
\node[st,below=of feat] (comp) {Compute six signals; calibrate image;\\grade chronology};
\node[st,below=of comp] (norm) {Drop missing signals, renormalise;\\compute evidence coverage};
\node[dec,below=of norm] (min) {Score $\geq$ 0.40?};
\node[st,right=14mm of min] (drop) {Retire any stored pair\\(no current lead)};
\node[st,below=of min] (store) {Store Match: tier capped by\\coverage and chronology};
\node[dec,below=of store] (tier) {Strong or\\moderate?};
\node[st,right=14mm of tier] (quiet) {Visible on report page,\\no notification};
\node[st,below=of tier] (notify) {Notify both reporters\\(once per pair)};
\draw[arr] (s)--(cand); \draw[arr] (cand)--(feat);
\draw[arr] (feat)--(comp); \draw[arr] (comp)--(norm); \draw[arr] (norm)--(min);
\draw[arr] (min)--node[lab,above]{no}(drop); \draw[arr] (min)--node[lab,right]{yes}(store);
\draw[arr] (store)--(tier);
\draw[arr] (tier)--node[lab,above]{weak}(quiet); \draw[arr] (tier)--node[lab,right]{yes}(notify);
\end{tikzpicture}
\caption{Matching algorithm flowchart}
\end{figure}

\subsection{Stable Text Model}\label{sec:text}
\begin{itemize}
  \item \textbf{Problem in \code{synopsis-v2}:} TF--IDF was fitted on the query report and its current candidate pool, so a pair's text score changed whenever an unrelated report was added, and differed depending on which side started matching.
  \item \textbf{Fix:} a fixed vocabulary compiled from a curated lexicon into a versioned release (\code{matching/text\_model/lexical-v1.json}). Each report becomes an L2-normalised map of term $\rightarrow$ class weight $\times (1 + \ln \mathit{tf})$ over its \emph{title and description}; similarity is the cosine of the two maps.
  \item \textbf{Normalisation:} synonyms and paraphrases (``mobile phone'' = ``cellphone''; ``cracked screen'' $\approx$ ``damaged display''), an object hierarchy, brand aliases, spelling variants, plurals and campus terms. Identifiers such as model numbers are kept exactly and weighted highest.
  \item \textbf{Class weights instead of IDF:} identifier 2.5, brand and object 2.0, feature 1.5, material 1.2, colour 1.0, descriptor 0.6, place 0.5. A corpus IDF would need either the live pool (the source of instability) or the evaluation data (leakage). Sentence-BERT embeddings \cite{sbert} remain an optional, separately versioned text profile.
  \item \textbf{Guarantees, each tested:} the same pair scores the same in both directions and whatever else is stored; releases are immutable; changing the release changes the engine version and invalidates cached features.
\end{itemize}
\begin{table}[H]
\centering\small
\caption{Text similarity examples under the lexical release}
\begin{tabularx}{\textwidth}{X X c}
\toprule
Report A & Report B & Text score\\
\midrule
Black mobile phone with cracked screen & Dark cellphone with damaged display & 0.95\\
Casio fx-991ES calculator & Casio FX991ES & 0.85\\
Casio fx-991ES calculator & Casio fx-82MS & 0.56\\
iPhone 13 with blue case & Blue Apple phone found on table & 0.77 (was 0.52)\\
\bottomrule
\end{tabularx}
\end{table}

\subsection{Image Similarity Calibration}\label{sec:calib}
\begin{itemize}
  \item \textbf{Problem:} pretrained image features never place unrelated photographs at zero. ResNet-50's pooled features \cite{resnet} are all non-negative and CLIP's embeddings \cite{clip} share a large common component.
  \item \textbf{Effect before the fix:} with ResNet-50, two unrelated photographs typically added about $0.275 \times 0.30 \approx 0.08$ to the total score.
  \item \textbf{Fix (introduced in \code{synopsis-v2}, kept in v3):} rescale the cosine above each model's measured chance level $c$ and clamp to $[0,1]$:
\end{itemize}
\begin{equation}
s_{\text{image}} \;=\; \min\!\left(1,\ \max\!\left(0,\ \frac{\cos(\mathbf{v}_a,\mathbf{v}_b) - c}{1 - c}\right)\right).
\end{equation}
\begin{itemize}
  \item \textbf{Result:} unrelated photos score about 0, identical photos still score 1, and candidate order is unchanged because the mapping is monotonic.
  \item \textbf{Chance levels:} the median cosine between photographs of \emph{different} items on the 120 negative pilot pairs, stored in \code{matching/profile.json}.
\end{itemize}

\begin{figure}[!htbp]
\centering
\begin{minipage}[c]{0.52\textwidth}
\centering
\begin{tikzpicture}[x=6.2cm,y=4.2cm]
\draw[ink!30] (0,0) grid[xstep=0.25,ystep=0.25] (1,1);
\draw[-{Latex[length=2mm]}] (0,0) -- (1.06,0) node[right,font=\scriptsize]{raw};
\draw[-{Latex[length=2mm]}] (0,0) -- (0,1.08) node[above,font=\scriptsize]{calibrated};
\foreach \v in {0,0.25,0.5,0.75,1} {\node[font=\tiny,below] at (\v,0) {\v}; \node[font=\tiny,left] at (0,\v) {\v};}
\draw[ink!50,dashed] (0,0) -- (1,1);
\draw[llaccent,very thick] (0,0) -- (0.275,0) -- (1,1);
\draw[teal,very thick] (0,0) -- (0.568,0) -- (1,1);
\draw[llfound,very thick,densely dotted] (0,0) -- (0.574,0) -- (1,1);
\fill[llaccentink] (0.857,0.803) circle (0.9mm);
\draw[llaccentink,thin] (0.857,0.803) -- (0.66,0.93); \node[font=\tiny,anchor=east,fill=paper,inner sep=1pt] at (0.66,0.93) {ID-card pair: 0.857 $\rightarrow$ 0.803};
\end{tikzpicture}
\end{minipage}\hfill
\begin{minipage}[c]{0.44\textwidth}
\small
\begin{tabular}{@{}llc@{}}
\toprule
& Model & Chance $c$\\
\midrule
\tikz\draw[llaccent,very thick] (0,0)--(0.5,0); & ResNet-50 (deployed) & 0.275\\
\tikz\draw[teal,very thick] (0,0)--(0.5,0); & CLIP (optional) & 0.568\\
\tikz\draw[llfound,very thick,densely dotted] (0,0)--(0.5,0); & Colour histogram & 0.574\\
\tikz\draw[ink!50,dashed] (0,0)--(0.5,0); & No calibration & --\\
\bottomrule
\end{tabular}
\end{minipage}
\caption{Calibration maps each model's raw cosine to a 0--1 score above its chance level}
\end{figure}

\begin{lstlisting}[style=lostlink,caption={Calibrated image similarity (\code{matching/engine.py}, abridged)}]
def calibrated_image_similarity(vec_a, vec_b, kind=None):
    raw = vector_similarity(vec_a, vec_b)          # cosine of feature vectors
    kind = kind or configuration()['image']        # resnet | clip | histogram
    chance = settings.MATCHING['IMAGE_CHANCE'].get(kind, 0.0)
    return min(1.0, max(0.0, (raw - chance) / (1.0 - chance)))
\end{lstlisting}

\subsection{Brand and Colour Matching}\label{sec:attr}
\begin{itemize}
  \item \textbf{Brands:} compared on letters and digits with an alias list, so ``Hydro Flask'' = ``hydroflask'' and ``HP'' = ``Hewlett-Packard''.
  \item \textbf{Colours as sets:} each description becomes a set of colour families (navy, sky and royal are blue; maroon is red; gray is grey); two sets are compared with a soft F1 in which neighbouring families earn partial credit.
  \item \textbf{Partial is not full:} sharing one colour no longer counts as complete agreement.
\end{itemize}
\begin{table}[H]
\centering\small
\caption{Colour comparison examples}
\begin{tabularx}{\textwidth}{X X c l}
\toprule
Colour A & Colour B & Score & Grade\\
\midrule
navy blue & blue & 1.00 & full\\
dark red & maroon & 1.00 & full\\
red and blue & red & 0.67 & partial (was 1.00)\\
light blue & navy & 0.75 & partial\\
grey & silver & 0.50 & neighbouring\\
red & blue & 0.00 & contradictory\\
\bottomrule
\end{tabularx}
\end{table}

\subsection{Chronology and Evidence Coverage}\label{sec:chrono}
\begin{itemize}
  \item \textbf{Problem:} the date signal ignored the reported time, and any found report dated more than a day before the loss was silently discarded, even when the mistake was probably a mis-typed date.
  \item \textbf{Fix:} date and optional time are compared together (across midnight too). Missing time is treated separately from a missing date, and nothing is excluded outright: the pair is graded, penalised and shown to moderators with a warning.
\end{itemize}
\begin{table}[H]
\centering\small
\caption{Chronology grades (found compared with lost)}
\begin{tabularx}{\textwidth}{P{24mm}X P{40mm}}
\toprule
Grade & When & Effect\\
\midrule
Plausible & Found at or after the loss & Normal proximity score\\
Approximate & Up to 1 day earlier by date, or up to 3 hours earlier by time & Date score $\times 0.9$\\
Suspicious & Earlier than that, within 7 days (e.g.\ lost 18:00, found 09:00 the same day) & $\times 0.5$; tier capped at moderate; warning\\
Contradictory & More than 7 days earlier & Date score 0; tier capped at weak; warning\\
\bottomrule
\end{tabularx}
\end{table}
\begin{itemize}
  \item \textbf{Missing versus contradictory:} a component is marked \emph{missing} (no photo on one report, no brand or colour, photo not processed, different image models) or \emph{conflict} (evidence present and disagreeing). Missing evidence is dropped; contradicting evidence counts against the pair.
  \item \textbf{Coverage on screen:} every match card shows the coverage beside the score and, for each signal, its base weight, effective weight after renormalisation (e.g.\ 30\,\% $\rightarrow$ 46\,\%) and points.
\end{itemize}

\subsection{Operational Traceability}
\begin{itemize}
  \item Every stage of the pipeline leaves evidence that a moderator or a test can inspect.
  \item Missing evidence lowers the number of signals; it never produces a guessed score.
\end{itemize}
\begin{table}[H]
\centering\small
\caption{Matching pipeline controls and recorded evidence}
\begin{tabularx}{\textwidth}{P{28mm}X X}
\toprule
Stage & Control & Recorded evidence\\
\midrule
Input validation & Reject invalid formats, oversized images and impossible dates & Form errors; database unchanged\\
Feature extraction & Cache vectors against a fingerprint of text, photo, model and text release & \code{ItemFeatures}: success / pending / missing / failed / retrying, model actually used\\
Score calculation & Calibrate image cosine, grade chronology, drop missing signals, renormalise & Component scores and status, coverage, chronology grade\\
Publication & Match only approved, active, unreserved opposite-type reports & Eligibility and engine version\\
Lifecycle & Retire (not delete) pairs whose report closed, changed or fell below 0.40 & Retired reason and time; history kept for claims\\
Taxonomy change & Rescore pairs affected by a changed related category or adjacent zone & Taxonomy fingerprint on each match\\
Claim review & Keep similarity separate from ownership approval & Questions, answers, decision, handover\\
\bottomrule
\end{tabularx}
\end{table}

\section{Notifications}
\begin{itemize}
  \item Fifteen events create in-app notifications; each one names who receives it. No notification ever contains a handover code or private evidence.
\end{itemize}
\begin{table}[H]
\centering\small
\caption{Notification events and recipients}
\begin{tabularx}{\textwidth}{X P{38mm} P{18mm}}
\toprule
Event & Recipient(s) & Kind\\
\midrule
Report approved (with number of candidates found) & Reporter & Report\\
Report rejected, or removed after a flag & Reporter & Report\\
New strong or moderate similarity match & Both reporters, once per pair & Match\\
Someone claims or offers to return an item & Reporter of the item & Claim\\
Verification question asked / answers submitted & Claimant / reporter & Claim\\
Claim approved: handover pending (how the code works) & Owner and the person handing over & Claim\\
Handover code created (the code itself is never sent) & Person handing over & Claim\\
Handover code locked after failed attempts & Owner & Claim\\
Handover completed & Claimant and reporter & Claim\\
Claim rejected, closed or revoked & Claimant & Claim\\
Claim withdrawn / reopened for reconsideration & Reporter of the item & Claim\\
Dispute raised / resolved & The other participants / both & Claim\\
Role changed by an administrator & That user & System\\
\bottomrule
\end{tabularx}
\end{table}
\begin{itemize}
  \item Notifications are private to their recipient and de-duplicated: a repeated unread event is reused, not duplicated.
  \item The masthead badge polls an unread-count endpoint; the inbox can mark one or all as read.
\end{itemize}

\section{Ownership Workflow and Handover}\label{sec:handover}
\begin{itemize}
  \item \textbf{Four concerns, kept apart:} the claim's verification status, the ownership \emph{reservation} (an approved claim, unique per item in the database), disputes (a separate record) and the physical handover (a code session). The central rule: claim approved $\neq$ item returned.
  \item \textbf{Independent decision:} approval needs every private question answered and a four-part evidence checklist (statement, private verification details, answers, supporting evidence); it is refused when any check contradicts the claim or none supports it, and the reason is recorded.
  \item \textbf{Disputes:} raising one keeps the claim's status and reservation and suspends any active code. An independent moderator dismisses it (the claim continues), revokes an approval (the item is released and matching resumes) or rejects an open claim. Reassignment is therefore always two audited steps: revoke, then approve another claim.
  \item \textbf{Reconsideration:} a rejected claimant may resubmit with new information and evidence (at most two requests, three per day); every statement is kept as an immutable revision, and it is refused while someone else holds the reservation.
\end{itemize}
\begin{figure}[!htbp]
\centering
\begin{tikzpicture}[scale=0.86,every node/.append style={transform shape},node distance=4mm and 9mm,
  st/.style={box,minimum width=44mm,minimum height=8mm}]
\node[term] (a) {Claim approved};
\node[st,below=of a] (meet) {Owner and the other person meet};
\node[st,below=of meet] (gen) {Owner creates a code:\\\code{secrets} 6 digits, 10-minute expiry};
\node[st,below=of gen] (store) {Server stores salted HMAC only;\\older active code superseded};
\node[st,below=of store] (tell) {Owner checks the item, then\\tells the code};
\node[st,below=of tell] (enter) {Other person enters it};
\node[dec,below=of enter] (ok) {Valid, unexpired,\\no dispute?};
\node[st,right=of ok] (fail) {Count attempt;\\lock after 5};
\node[st,below=of ok] (done) {Claim and report returned;\\counterpart report closed};
\node[term,below=of done] (end) {Both notified};
\node[st,right=of gen,fill=rust!8] (disp) {Dispute raised:\\code suspended,\\reservation kept};
\draw[arr] (a)--(meet); \draw[arr] (meet)--(gen); \draw[arr] (gen)--(store); \draw[arr] (store)--(tell);
\draw[arr] (tell)--(enter); \draw[arr] (enter)--(ok); \draw[arr] (ok)--node[lab,above]{no}(fail);
\draw[arr] (ok)--node[lab,right]{yes}(done); \draw[arr] (done)--(end);
\draw[arr,dashed] (disp.west)--(gen.east);
\end{tikzpicture}
\caption{One-time-code handover}
\end{figure}
\begin{table}[H]
\centering\small
\caption{Handover code safeguards}
\begin{tabularx}{\textwidth}{P{38mm}X}
\toprule
Property & Implementation\\
\midrule
Who holds the code & Always the owner: the claimant on a found report, the reporter on a lost report\\
Generation & \code{secrets.randbelow(10**6)}, shown once on a no-store page\\
Storage & HMAC-SHA-256 with a per-session salt; plaintext never stored, logged or notified\\
Lifetime and use & 10 minutes; single use; a new code supersedes the old one (one active code per claim, database constraint)\\
Guessing & 5 wrong entries lock the code; 10 verifications per user per 10 minutes; 5 codes per claim per hour\\
Preconditions & Approved claim, no open dispute, open report; checked with the item and claim rows locked\\
Moderator role & Sees status, attempts and times, never the code; may record an \emph{exception} only with a reason and written evidence, logged separately\\
\bottomrule
\end{tabularx}
\end{table}
\begin{lstlisting}[style=lostlink,caption={Verifying a handover code (\code{claims/services.py}, abridged)}]
def verify_handover_code(claim, item, user, code):
    if user.id != claim.provider.id: raise ClaimError('Only the person handing over')
    if handover_blocker(claim, item): raise ClaimError(...)   # approved, no dispute
    session = claim.handover_sessions.select_for_update().get(status='active')
    if session.is_expired: return expire(session)
    if not constant_time_compare(_hash_code(session.salt, code), session.otp_hash):
        session.failed_attempts += 1                       # persisted, not rolled back
        if session.failed_attempts >= MAX_ATTEMPTS: session.status = 'locked'
        ...
    complete_handover(claim, item, user, method='otp')     # claim + report returned
\end{lstlisting}

\section{Security Controls}
\begin{itemize}
  \item Each threat below is handled on the server, so hiding a button is never the only protection.
\end{itemize}
\begin{table}[H]
\centering\small
\caption{Threats and the controls that address them}
\begin{tabularx}{\textwidth}{P{38mm}X}
\toprule
Threat & Control in LostLink\\
\midrule
Cross-site request forgery & CSRF token on every POST (sign-out included); \code{SameSite=Lax} cookies\\
Session theft & \code{HttpOnly} session cookie; \code{Secure} cookies, HTTPS redirect and HSTS in production\\
Clickjacking, sniffing, referrers & \code{X-Frame-Options: DENY}; \code{nosniff}; \code{Referrer-Policy: same-origin}\\
Non-members and fake accounts & Expiring single-use email verification; optional campus domain list\\
Password guessing, enumeration & Rate limits on sign-in, registration, resend and reset; identical responses for known and unknown addresses\\
Privilege escalation & Role never accepted from forms; admin role fields read-only for non-superusers\\
Insecure direct object access & Ownership and participation checks on reports, claims, answers and evidence\\
Photo exposure and metadata & Deny-by-default media; uploads re-encoded without EXIF/GPS; sensitive photos blurred for others\\
Malicious uploads & Signature check, full decode, size and pixel limits, polyglot rejection, random names\\
Verification secrets in public & Identifier guard on public text; private details and expected answers hidden\\
Handover fraud & Owner-held hashed one-time code; attempt lock; disputes suspend it\\
Tampering with history & Append-only audit records (updates and deletes refused); read-only admin\\
CSV formula injection & Cells starting with \code{=}, \code{+}, \code{-}, \code{@}, tab or line breaks are escaped\\
Unsafe production settings & Startup refuses short secret keys, wildcard hosts and non-https CSRF origins\\
Conflict of interest & Nobody reviews their own report, decides their own case or resolves a dispute they raised\\
\bottomrule
\end{tabularx}
\end{table}

\section{Workflow Integrity Rules}
Two logic audits of the full recovery journey found defects where the stored state could disagree with reality or where the system claimed more than it knew. Each was fixed and is guarded by a regression test.
\begin{table}[H]
\centering\small
\caption{Workflow defects found in the first logic audit and the rule now enforced}
\begin{tabularx}{\textwidth}{P{4mm}X X}
\toprule
\# & Defect before the audit & Rule now enforced\\
\midrule
1 & A moderator could decide a claim on an item they had reported. & The finder can never decide the claim (HTTP 403).\\
2 & A confirmed return left the claimant's own lost report open and matching. & Claimant's report marked returned; its matches and claims closed.\\
3 & Removing a flagged report left its claims open. & Rejecting the report closes unfinished claims and notifies claimants.\\
4 & Withdrawing a report left its flags in the queue. & Withdrawal resolves the report's open flags.\\
5 & An impossible analytics date (30 February) crashed the page. & Invalid dates are ignored; the full period is shown.\\
6 & A non-numeric user ID caused a server error. & The ID is validated before any lookup.\\
7 & Protected pages dropped the requested address at sign-in. & The original path is kept in \code{?next=}.\\
8 & Two accounts could share an email differing only in case. & Emails are lower-cased and checked for uniqueness.\\
\bottomrule
\end{tabularx}
\end{table}
\begin{itemize}
  \item \textbf{Live-only statistics:} analytics had counted 14 stale matches out of 20 because they pointed at already-returned reports; dashboards now count only live matches.
  \item \textbf{Evidence kept:} match rows are retired rather than deleted, so closed claims keep their evidence.
  \item \textbf{Chart fix:} the weekly-reports chart used two undefined colour tokens and drew invisible bars; it now uses the theme's lost and found colours (commit \code{62334a7}).
\end{itemize}

\begin{table}[!tbp]
\centering\small
\caption{Defects found in the second (remediation) audit and the rule now enforced}
\begin{tabularx}{\textwidth}{P{4mm}X X}
\toprule
\# & Defect before the remediation & Rule now enforced\\
\midrule
1 & Disputing an approved claim changed its status and released the reservation, letting a rival claim proceed. & Disputes are separate records; the reservation stays until an independent moderator revokes it.\\
2 & A moderator could approve or reject their own report. & Self-review refused on every path; own reports listed separately.\\
3 & A moderator's single click marked an item returned. & Only the owner's one-time code, or a documented exception, returns an item.\\
4 & A rejected claimant was sent back to the closed claim with no way forward. & Limited, audited reconsideration with statement history.\\
5 & ``Claim this item'' showed on items already reserved, then failed. & Buttons follow server-side eligibility; ``Handover pending'' shown instead.\\
6 & A pair's text score changed when unrelated reports were added. & Fixed, versioned text model; tested invariants.\\
7 & ``red and blue'' matched ``red'' fully; found-before-lost pairs were silently hidden. & Colour sets; graded chronology with warnings.\\
8 & Anyone could register any email and act at once; no password recovery. & Email verification, domain policy, password reset, rate limits.\\
9 & Uploaded photos kept EXIF/GPS data; ID cards were shown to everyone. & Re-encoding without metadata; sensitive photos blurred.\\
10 & A fresh production database had no categories or zones. & \code{seed\_taxonomy} initialises them without demo accounts.\\
\bottomrule
\end{tabularx}
\end{table}

\section{API Integration}
\begin{itemize}
  \item No public third-party API is needed for the core workflow.
  \item Internal JSON endpoints serve search results, unread counts and read-state updates; none return contact details, answers or full-resolution sensitive photos.
  \item Pretrained weights are an optional download at setup; if inference fails, the failure is recorded in \code{ItemFeatures}, matching continues with the remaining evidence, and a baseline model is never presented as a pretrained one.
  \item Outgoing email goes through Django's SMTP backend configured from the environment; with no mail server in production the app logs each undelivered message and the system check warns.
\end{itemize}

% ===========================================================================
\chapter{Module Description}
\textbf{In this chapter:}
\begin{itemize}[beginpenalty=10000]
  \item how the six Django apps depend on each other;
  \item what each module is responsible for;
  \item the size and test coverage of every module.
\end{itemize}

\begin{figure}[!htbp]
\centering
\begin{tikzpicture}[
  app/.style={draw=ink,rounded corners=1.5mm,fill=pale,minimum width=30mm,minimum height=9mm,font=\small\ttfamily},
  dep/.style={-{Latex[length=2mm]},thick,color=ink!60}]
\node[app] (dash) at (5,0) {dashboard};
\node[app] (claims) at (1.2,-1.7) {claims};
\node[app] (match) at (8.8,-1.7) {matching};
\node[app] (items) at (5,-3.4) {items};
\node[app] (acc) at (0.6,-3.4) {accounts};
\node[app] (notes) at (5,-5.1) {notifications};
\draw[dep] (dash) -- (claims);
\draw[dep] (dash) -- (match);
\draw[dep] (claims) -- (match);
\draw[dep] (claims) -- (items);
\draw[dep] (claims.west) -| (acc.north);
\draw[dep] (match) -- (items);
\draw[dep] (items) -- (acc);
\draw[dep] (claims.south) .. controls +(0,-2.4) and +(-2.2,0) .. (notes.west);
\draw[dep] (match.south) .. controls +(0,-2.4) and +(2.2,0) .. (notes.east);
\end{tikzpicture}
\caption[Module dependencies]{Module dependencies (arrow = ``uses''); \code{items} also triggers \code{matching} when a report changes.}
\end{figure}

\section{Module 1: User Management (\code{accounts})}
\begin{itemize}
  \item Registration with email verification, sign-in, POST sign-out, password reset, profile editing and the custom role model.
  \item Role and verification decorators that preserve the requested page across sign-in; a shared cache-backed rate limiter.
  \item Initials and display names give a consistent identity without profile photographs.
\end{itemize}

\section{Module 2: Core Functionality (\code{items}, \code{matching}, \code{claims})}
\begin{itemize}
  \item \code{items}: the validated report lifecycle (\code{items/lifecycle.py}), search, browse, the landing page, upload sanitising and the deny-by-default media view.
  \item \code{matching}: the versioned text model, feature status, calibration, chronology, colour sets, coverage, scoring, retirement and evaluation commands.
  \item \code{claims}: the service layer (\code{claims/services.py}) for verification, checklist decisions, reservation, disputes, reconsideration, private evidence and the one-time-code handover.
\end{itemize}

\section{Module 3: Data Management}
\begin{itemize}
  \item Feature vectors are cached against fingerprints of the text, photo, model and text release; matches carry their engine version and taxonomy fingerprint.
  \item Editing, withdrawing or reserving a report retires its matches; history stays for claims.
  \item Evidence files are purged after the retention period; audit records are append-only.
  \item \code{rebuild\_matches} recomputes every pair after an engine change without touching reports or claims.
  \item Pagination and filters keep every list operation bounded.
\end{itemize}

\section{Module 4: Admin Functionality (\code{dashboard})}
\begin{itemize}
  \item Moderators: review queue for reports (own reports listed separately), claims, disputes, handovers awaiting a code, flags and sensitive-photo review; decision log with previous and new state.
  \item Administrators: roles and account status, categories and zones, analytics, CSV and PDF exports.
\end{itemize}

\section{Module 5: Additional Features (\code{notifications} and UI)}
\begin{itemize}
  \item Persistent, de-duplicated in-app notifications with an unread badge.
  \item Light/Dark/System themes, a floating phone tab bar and reduced-motion handling.
  \item Shareable search URLs, development cache busting, CSV formula-injection protection, PDF summaries and score-drift snapshots.
\end{itemize}

\begin{table}[H]
\centering\small
\renewcommand{\arraystretch}{1.05}
\caption{Module responsibilities, size and test coverage}
\begin{tabularx}{\textwidth}{P{22mm}P{34mm}X rr}
\toprule
Module & Responsibility & Principal output & Lines & Tests\\
\midrule
accounts & Identity and access & Verified user, role predicates, rate limits & 1{,}818 & 70\\
items & Report lifecycle and discovery & Approved, searchable reports; safe media & 3{,}279 & 70\\
matching & Text model, features, scoring & Versioned, explained, stable match & 5{,}771 & 120\\
claims & Ownership verification & Reserved claim; code-confirmed handover & 3{,}017 & 67\\
notifications & User attention & Private, de-duplicated events & 261 & 6\\
dashboard & Governance and insights & Decisions, analytics, exports & 1{,}182 & 23\\
\midrule
\multicolumn{3}{l}{Total (8 configuration tests included)} & 15{,}328 & 364\\
\bottomrule
\end{tabularx}
\par\vspace{2pt}{\footnotesize Python lines include tests and exclude migrations. The interface adds 56 templates, 1{,}409 lines of CSS and 549 of JavaScript.}
\end{table}

% ===========================================================================
\chapter{User Interface}
\textbf{In this chapter:}
\begin{itemize}[beginpenalty=10000]
  \item the site map and the design tokens behind every page;
  \item a walkthrough of each screen, from the landing page to analytics;
  \item dark theme and phone layouts.
\end{itemize}
\begin{itemize}
  \item Screenshots were captured from the running application (commit \code{3eda809}, engine \code{synopsis-v3}) on a copy of the local demonstration database.
  \item Desktop captures use a 1440-pixel width; phone captures use 390 pixels.
  \item Demo reports use representative photographs; the ID-card pair uses sample photographs prepared for this report.
\end{itemize}

\section{Navigation and Design System}
\begin{itemize}
  \item \textbf{Role-based navigation:} each role sees its own menu; higher roles include the pages of lower ones.
  \item \textbf{Design tokens:} one set of CSS custom properties drives colours, type, spacing and motion in both themes.
  \item \textbf{Consistent markers:} vermilion always means \emph{lost}, warm grey always means \emph{found}.
\end{itemize}
\begin{figure}[!htbp]
\centering
\begin{tikzpicture}[scale=0.8,every node/.append style={transform shape},
  pg/.style={draw=ink,rounded corners=1mm,fill=paper,font=\scriptsize,minimum width=29mm,minimum height=6.2mm,anchor=west},
  grp/.style={draw=ink,rounded corners=1mm,fill=pale,font=\scriptsize\bfseries,minimum width=30mm,minimum height=7mm},
  e/.style={draw=ink!60}]
\node[grp] (root) at (9.3,0) {Landing page};
\foreach \g/\x/\t in {pub/2/Public,stu/7.1/Student,mod/12.2/Moderator,adm/17.3/Administrator}
  \node[grp] (\g) at (\x,-1.5) {\t};
\foreach \g in {pub,stu,mod,adm} \draw[e] (root.south) -- ++(0,-0.4) -| (\g.north);
\foreach \i/\t in {1/Sign in,2/Register and verify,3/Password reset,4/Search,5/Browse lost and found}
  {\node[pg] (p\i) at (0.9,-1.65-0.8*\i) {\t}; \draw[e] (0.6,-1.85) |- (p\i.west);}
\foreach \i/\t in {1/Home,2/Report an item,3/Report detail,4/My reports,5/Claims,6/Notifications,7/Profile}
  {\node[pg] (s\i) at (6.0,-1.65-0.8*\i) {\t}; \draw[e] (5.7,-1.85) |- (s\i.west);}
\foreach \i/\t in {1/Review queue,2/Checklist decisions,3/Disputes and flags,4/Handover status,5/Decision log}
  {\node[pg] (m\i) at (11.1,-1.65-0.8*\i) {\t}; \draw[e] (10.8,-1.85) |- (m\i.west);}
\foreach \i/\t in {1/Analytics,2/CSV and PDF exports,3/People and roles,4/Taxonomy}
  {\node[pg] (a\i) at (16.2,-1.65-0.8*\i) {\t}; \draw[e] (15.9,-1.85) |- (a\i.west);}
\end{tikzpicture}
\caption[Site map]{Site map. Moderators also see student pages; administrators also see moderator pages.}
\end{figure}

\begin{table}[H]
\centering\small
\caption{Core design tokens of the Black + Vermilion theme (light mode)}
\begin{tabularx}{\textwidth}{P{8mm}P{30mm}P{20mm}X}
\toprule
& Token & Value & Used for\\
\midrule
\swatch{llpage} & \code{--page} & \#F5F3F0 & Warm off-white page background\\
\swatch{white} & \code{--surface} & \#FFFFFF & Cards and panels\\
\swatch{lltext} & \code{--text} & \#211F1E & Body text and headings\\
\swatch{llaccent} & \code{--accent}, \code{--lost-dot} & \#D6533B & Brand, focus rings and lost markers\\
\swatch{llaccentink} & \code{--accent-ink} & \#B83E2B & Accent-coloured text (meets contrast)\\
\swatch{llfound} & \code{--found-dot} & \#827B75 & Found markers and neutral chart bars\\
\swatch{llsoft} & \code{--accent-soft} & \#FFF1ED & Action banners and tints\\
\bottomrule
\end{tabularx}
\end{table}

\section{Landing Page}
\begin{itemize}
  \item Typed headline and an orbit of belongings circling the brand mark.
  \item The two photo cards show the \textbf{most similar live pair on campus}, here the lost and found student ID cards, linked by their similarity score of 92\%.
  \item Calls to action lead to reporting or browsing; live counts show open reports, zones and strong or moderate matches.
  \item Photographs marked sensitive are never used on the landing page.
\end{itemize}
\begin{figure}[!htbp]
\centering
\shotframe[0.97\textwidth]{landing.jpg}
\caption[Landing page]{Landing page. The Lost and Found cards in the orbit are the most similar live pair, joined by its real 92\% similarity score.}
\end{figure}

\section{Sign-in and Registration}
\begin{itemize}
  \item Two-panel layout: the form, and a four-step summary of how recovery works.
  \item Registration always creates a Student, signs nobody in and emails a verification link; the password checklist updates while typing.
  \item ``Forgot your password?'' starts an emailed reset; responses never reveal whether an address is registered.
  \item Demo credentials appear only in development mode.
\end{itemize}
\begin{figure}[!htbp]
\centering
\begin{subfigure}[t]{0.49\textwidth}
\shotframe{login.jpg}
\caption{Sign-in (development mode)}
\end{subfigure}\hfill
\begin{subfigure}[t]{0.49\textwidth}
\shotframe{register.jpg}
\caption{Account registration}
\end{subfigure}
\caption{Public entry points}
\end{figure}

\section{Student Home}
\begin{itemize}
  \item Leads with the number of similar reports to check and the most similar one, drawn as two discs whose overlap equals the score.
  \item Pending actions (here, a verification question) appear as a banner.
  \item Ranked similar reports, claims in progress and recent notifications follow.
\end{itemize}
\begin{figure}[!htbp]
\centering
\shotframe[0.97\textwidth]{student-home.jpg}
\caption{Student home with the strongest match, a pending action and ranked matches}
\end{figure}

\section{Reporting and Search}
\begin{itemize}
  \item \textbf{Report form:} what happened, what the item is, where and when, optional brand, colour and photo, with a live preview card.
  \item \textbf{Private verification details:} a separate box, visible only to the reporter and moderators, for what the real owner would know; ID, serial and phone numbers typed into the public description are refused with a pointer to this box.
  \item \textbf{Search:} keyword, type, category and zone in one panel; results update through debounced AJAX and the URL always reflects the filters.
\end{itemize}
\begin{figure}[!htbp]
\centering
\begin{subfigure}[t]{0.49\textwidth}
\shotframe{report-form.jpg}
\caption{Report form with live preview}
\end{subfigure}\hfill
\begin{subfigure}[t]{0.49\textwidth}
\shotframe{report-private.jpg}
\caption{Public description and the private details box}
\end{subfigure}
\caption{Reporting}
\end{figure}

\section{Report Detail and Match Explanation}
\begin{itemize}
  \item Photograph, description and structured facts, including which neighbouring zones earn partial credit.
  \item Each candidate shows its similarity score, heuristic tier and evidence coverage, with the reminder that a score is not a probability or proof of ownership.
  \item \emph{Why this score} lists all six signals with their status, base and effective weight and points; the points add up to the total, and missing signals say why.
\end{itemize}
\begin{figure}[!htbp]
\centering
\shotframe[0.97\textwidth]{item-detail.jpg}
\caption{Lost report for the JIIT student ID card}
\end{figure}
\begin{figure}[!htbp]
\centering
\shotframe[0.97\textwidth]{match-explain.jpg}
\caption[Per-signal explanation of the ID-card match]{Per-signal explanation of the 92\% ID-card match (worked through in Section~\ref{sec:worked})}
\label{fig:explain}
\end{figure}

\section{Claims, Verification and Handover}
\begin{itemize}
  \item The moderator records an \textbf{evidence checklist} and a reason; the reporter's private details and the expected answers are shown only to staff and the reporter.
  \item After approval the owner sees \textbf{Handover pending}: the item is reserved and contacts are shared.
  \item At the meeting the owner creates a \textbf{one-time code}, shown once; the person handing the item over enters it; moderators see only its status.
\end{itemize}
\begin{figure}[!htbp]
\centering
\begin{subfigure}[t]{0.49\textwidth}
\shotframe{claim-review.jpg}
\caption{Evidence checklist for an independent moderator}
\end{subfigure}\hfill
\begin{subfigure}[t]{0.49\textwidth}
\shotframe{handover-owner.jpg}
\caption{Owner's claim page: approved, handover pending}
\end{subfigure}
\caption{Claim decision and reservation}
\end{figure}
\begin{figure}[!htbp]
\centering
\begin{subfigure}[t]{0.49\textwidth}
\shotframe{handover-code.jpg}
\caption{The owner's one-time code, shown once}
\end{subfigure}\hfill
\begin{minipage}[t]{0.49\textwidth}
\begin{subfigure}[t]{\textwidth}
\shotframe{handover-finder.jpg}
\caption{The finder enters the code}
\end{subfigure}\\[2mm]
\begin{subfigure}[t]{\textwidth}
\shotframe{handover-moderator.jpg}
\caption{Moderator view: status, never the code}
\end{subfigure}
\end{minipage}
\caption{One-time-code handover}
\end{figure}

\section{Privacy and Identity}
\begin{itemize}
  \item A photograph marked sensitive (here the ID card, marked by a moderator) is blurred for everyone except the reporter, staff and a verified owner.
  \item Accounts that have not confirmed their email can browse but not report, flag or claim until they open the link.
  \item My reports filters by state (active, returned, withdrawn or removed) and type; search shows only public reports.
\end{itemize}
\begin{figure}[!htbp]
\centering
\begin{subfigure}[t]{0.49\textwidth}
\shotframe{photo-blurred.jpg}
\caption{Sensitive photo seen by another student}
\end{subfigure}\hfill
\begin{subfigure}[t]{0.49\textwidth}
\shotframe{verify-pending.jpg}
\caption{Email verification pending}
\end{subfigure}
\\[3mm]
\begin{subfigure}[t]{0.49\textwidth}
\shotframe{search.jpg}
\caption{Search with combined filters}
\end{subfigure}\hfill
\begin{subfigure}[t]{0.49\textwidth}
\shotframe{my-reports.jpg}
\caption{My reports, filterable by state and type}
\end{subfigure}
\caption{Privacy, identity and discovery}
\end{figure}

\section{Staff Pages}
\begin{itemize}
  \item Review queue grouped into new reports, claims, disputes, handovers awaiting a code and flags, oldest first; a moderator's own reports are listed separately for someone else.
  \item Approve, reject and sensitive-photo controls; every decision goes to the append-only log with its previous and new state.
  \item Administrators manage people and roles; only they can grant staff access.
\end{itemize}
\begin{figure}[!htbp]
\centering
\begin{subfigure}[t]{0.49\textwidth}
\shotframe{moderation.jpg}
\caption{Moderator review queue}
\end{subfigure}\hfill
\begin{subfigure}[t]{0.49\textwidth}
\shotframe{users.jpg}
\caption{People and roles (administrator)}
\end{subfigure}
\caption{Moderation and administration}
\end{figure}

\section{Analytics}
\begin{itemize}
  \item Five headline figures and a weekly chart of new reports and recoveries (dated by the handover time).
  \item Breakdowns by status, claim outcome, similarity band, zone and category; recoveries are split into code-confirmed and exception handovers.
  \item An \emph{Operational data quality} panel that states plainly it does not measure accuracy; date filters and CSV or PDF export.
\end{itemize}
\begin{figure}[!htbp]
\centering
\shotframe[0.8\textwidth]{analytics.jpg}
\caption{Administrator analytics (engine \code{synopsis-v3})}
\end{figure}

\section{Dark Theme and Phone Layout}
\begin{itemize}
  \item Every page supports Light, Dark and System themes from the same tokens.
  \item On phones the masthead collapses and a floating tab bar offers Home, Search, Report, Claims and Alerts.
\end{itemize}
\begin{figure}[!htbp]
\centering
\begin{subfigure}[t]{0.49\textwidth}
\shotframe{landing-dark.jpg}
\caption{Landing page, dark theme}
\end{subfigure}\hfill
\begin{subfigure}[t]{0.49\textwidth}
\shotframe{home-dark.jpg}
\caption{Student home, dark theme}
\end{subfigure}
\\[3mm]
\begin{subfigure}[t]{0.25\textwidth}
\shotframe{mobile-home.jpg}
\caption{Phone home}
\end{subfigure}\hspace{6mm}
\begin{subfigure}[t]{0.25\textwidth}
\shotframe{mobile-item.jpg}
\caption{Phone report}
\end{subfigure}\hspace{6mm}
\begin{subfigure}[t]{0.25\textwidth}
\shotframe{mobile-search-dark.jpg}
\caption{Phone search, dark}
\end{subfigure}
\caption{Theme and responsive layouts}
\end{figure}

% ===========================================================================
\chapter{Testing}
\textbf{In this chapter:}
\begin{itemize}[beginpenalty=10000]
  \item the testing strategy, levels and tools;
  \item thirty representative test cases with their results;
  \item error handling and the verification evidence for this release.
\end{itemize}

\section{Testing Strategy}
\begin{itemize}
  \item \textbf{Risk-based:} permissions, state transitions and privacy get both positive and negative cases.
  \item \textbf{Layered:} many fast unit tests, fewer integration tests, and browser checks at the top.
  \item \textbf{Regression-first:} every audit defect was reproduced by a failing test before it was fixed.
  \item \textbf{Adversarial:} security tests try to bypass each rule directly (forged POSTs, wrong users, reused or guessed codes, tampered tokens), not only through the interface.
\end{itemize}
\begin{figure}[!htbp]
\centering
\begin{minipage}[c]{0.46\textwidth}
\centering
\begin{tikzpicture}[scale=0.95]
\fill[pale,draw=ink] (0,0) -- (7,0) -- (5.85,1.35) -- (1.15,1.35) -- cycle;
\fill[pale!60!white,draw=ink] (1.15,1.35) -- (5.85,1.35) -- (4.7,2.7) -- (2.3,2.7) -- cycle;
\fill[llsoft,draw=ink] (2.3,2.7) -- (4.7,2.7) -- (3.5,4.1) -- cycle;
\node[font=\scriptsize,align=center] at (3.5,0.66) {\textbf{Unit tests}\\text model, colour, chronology, roles};
\node[font=\scriptsize,align=center] at (3.5,2.0) {\textbf{Integration}\\workflows, handover, privacy};
\node[font=\scriptsize\bfseries] at (3.5,3.15) {UI};
\node[font=\scriptsize,align=center] at (3.5,4.6) {browser sweep and end-to-end journey (Playwright)};
\end{tikzpicture}
\end{minipage}\hfill
\begin{minipage}[c]{0.5\textwidth}
\centering
\begin{tikzpicture}[x=0.042cm,y=0.56cm]
\foreach \n/\v/\i in {accounts/70/0,items/70/1,matching/120/2,claims/67/3,dashboard/23/4,notifications/6/5,config/8/6}{
  \node[font=\scriptsize\ttfamily,anchor=east] at (-1,-\i) {\n};
  \fill[llaccent] (0,-\i-0.3) rectangle (\v,-\i+0.3);
  \node[font=\scriptsize,anchor=west] at (\v+1.2,-\i) {\v};}
\draw[ink!40] (0,0.5) -- (0,-6.5);
\node[font=\scriptsize\bfseries] at (55,0.95) {Automated tests per app (364)};
\end{tikzpicture}
\end{minipage}
\caption{Test pyramid and the distribution of the 364 automated tests}
\end{figure}

\begin{table}[H]
\centering\small
\caption{Test levels, tools and scope}
\begin{tabularx}{\textwidth}{P{26mm}P{38mm}X}
\toprule
Level & Tool & Scope\\
\midrule
Unit & \code{SimpleTestCase} & Text-model invariants, colour sets, chronology grades, calibration, coverage caps, tiers, form rules\\
Integration & \code{TestCase} & Verification, reset, moderation, matching, claims, disputes, appeals, handover codes, media, exports\\
Concurrency & \code{TransactionTestCase} & Constraints: one reserved claim per item, one open dispute and one active code per claim; double submissions\\
Regression & per-app regression modules & Every defect from both logic audits pinned by a test\\
First run & \code{smoke\_first\_run.py} & Empty database to code-confirmed handover in 11 steps\\
Browser sweep & Playwright \cite{playwright} & 4 roles $\times$ 24 routes $\times$ 3 widths, light and dark (318 checks)\\
End to end & Playwright script & Report $\rightarrow$ approve $\rightarrow$ match $\rightarrow$ claim $\rightarrow$ checklist $\rightarrow$ owner's code $\rightarrow$ returned\\
Mobile layout & Playwright script & 5 widths $\times$ 7 pages: content and last control clear of the tab bar (65 checks)\\
\bottomrule
\end{tabularx}
\end{table}

\section{Unit and Integration Testing}
\begin{itemize}
  \item \textbf{Unit:} isolated scoring logic with no database, plus model defaults and constraints.
  \item \textbf{Integration:} each workflow test runs against an isolated database and checks both the response and the stored state.
  \item \textbf{Regression:} every defect from both logic audits has a test that failed before its fix.
  \item \textbf{Handover (15 tests):} success in both lost and found flows, wrong, expired, reused and superseded codes, wrong users, codes before approval, disputes, double submission, attempt lock, rate limits and leakage of the code.
\end{itemize}

\section{Functional Testing}
\begin{itemize}
  \item \textbf{Browser sweep:} fails on HTTP errors, JavaScript errors, horizontal overflow and heading-structure faults.
  \item \textbf{End-to-end journey:} files a lost and a found report, approves both, checks the match, runs the private question flow, approves the claim with the checklist, has the owner create the code, checks the moderator never sees it, has the finder enter it and verifies that both reports are closed.
\end{itemize}

\section{Test Cases and Results}
\begin{itemize}
  \item Thirty representative cases: TC-01 to TC-12 cover the core workflow, TC-13 to TC-20 the calibration and first audit, TC-21 to TC-30 the remediation.
  \item Every case passes on the current code; each maps to at least one automated test.
\end{itemize}
\begin{longtable}{>{\bfseries}P{14mm}P{56mm}P{60mm}P{12mm}}
\caption{Representative test cases and results}\\
\toprule
ID & Scenario & Expected outcome & Result\\
\midrule
\endfirsthead
\toprule
ID & Scenario & Expected outcome & Result\\
\midrule
\endhead
TC-01 & Register with valid details & Unverified student created; link emailed; not signed in & Pass\\
TC-02 & Student opens a staff route & Access denied or redirected & Pass\\
TC-03 & Future-dated report or invalid image & Validation error; nothing saved & Pass\\
TC-04 & Approve a pending report & Published, matching runs, audit row stored & Pass\\
TC-05 & Search with combined filters & Only approved records matching every filter & Pass\\
TC-06 & Match a report without a photograph & Image dropped; weights renormalised & Pass\\
TC-07 & Found report dated 5 days before the loss & Stored as suspicious, capped at moderate, warning shown & Pass\\
TC-08 & Claim without an eligible relationship & Request rejected & Pass\\
TC-09 & Approve without answers, or with a contradicting check & Approval blocked & Pass\\
TC-10 & Approve two rival claims & Only one approved claim allowed & Pass\\
TC-11 & Export CSV with formula-like content & Cell escaped; private data omitted & Pass\\
TC-12 & Open private media as an unrelated user & Access denied & Pass\\
TC-13 & Compare photos of two different objects & Calibrated image score near zero & Pass\\
TC-14 & ``navy blue''/``blue''; ``red and blue''/``red'' & 1.00 full; 0.67 partial & Pass\\
TC-15 & Moderator decides a claim on own report & HTTP 403; claim unchanged & Pass\\
TC-16 & Finder enters the owner's valid code & Both reports returned; matches and claims closed & Pass\\
TC-17 & Reject a flagged report with open claims & Claims closed; claimants notified & Pass\\
TC-18 & Withdraw a report with an open flag & Flag resolved & Pass\\
TC-19 & Analytics with a date of 30 February & Full period shown, no error & Pass\\
TC-20 & Open a protected page while signed out & Sign-in keeps \code{?next=} & Pass\\
TC-21 & Dispute an approved claim; rival tries to claim & Reservation kept; rival refused; code suspended & Pass\\
TC-22 & Enter a wrong, expired or reused code & Refused; attempts counted; locked after 5 & Pass\\
TC-23 & Moderator creates or reads the code & Refused; only status visible & Pass\\
TC-24 & Unverified account starts a claim & Redirected to verification; nothing created & Pass\\
TC-25 & Moderator approves own report & Refused; report unchanged & Pass\\
TC-26 & Add unrelated reports, reverse direction & Pair score unchanged & Pass\\
TC-27 & Roll number typed in a public description & Refused; private details suggested & Pass\\
TC-28 & Upload a photo with GPS EXIF & Stored re-encoded with no metadata & Pass\\
TC-29 & Another student opens a sensitive photo & Blurred preview only & Pass\\
TC-30 & Third reconsideration request & Refused after the limit; history intact & Pass\\
\bottomrule
\end{longtable}

\section{Error Handling}
\begin{itemize}
  \item Invalid filter parameters and impossible dates are ignored instead of causing server errors.
  \item Model inference failures are recorded per report and retried in a controlled way; the match uses the remaining signals and says which were not compared.
  \item Custom 400, 403 (with the reason), 404, 500 and CSRF-failure pages give consistent recovery paths; rate-limited requests get a 429 page.
  \item Every state transition locks and re-reads the current rows, so stale, repeated or concurrent requests cannot corrupt a workflow; a failed code attempt is still recorded.
\end{itemize}

\begin{table}[H]
\centering\small
\caption{Repository verification evidence}
\begin{tabularx}{\textwidth}{P{46mm}X}
\toprule
Check & Recorded outcome\\
\midrule
\code{python manage.py check}, \code{makemigrations --check} & No issues; no missing migrations\\
\code{check --deploy} (production settings) & Only the expected warning that SMTP is not configured in the test environment\\
Automated suite & 364 passed on SQLite (commit \code{3eda809}); 1 opt-in pretrained test skipped; PostgreSQL 17 run by CI\\
First-run smoke test & 11 steps passed on an empty database (SMTP delivery not run)\\
Browser sweep & 318 page checks; no HTTP, overflow, JavaScript or heading faults\\
Mobile layout & 65 checks at 320--767\,px; content clear of the tab bar\\
End-to-end workflow & Full journey to a code-confirmed handover passes and closes both reports\\
Migration checks & Applied to a copy of the populated demo database; legacy disputes and accounts migrated, data kept\\
\bottomrule
\end{tabularx}
\end{table}

\begin{table}[H]
\centering\small
\caption{Testing concerns and assurance level}
\begin{tabularx}{\textwidth}{P{30mm}X P{26mm}}
\toprule
Concern & Representative assurance & Level\\
\midrule
Access control & Role gates, ownership, self-decision and unrelated-user denial & High\\
State integrity & Lifecycle, reservation, dispute, appeal and code-handover transitions & High\\
Privacy & Media access, sensitive photos, metadata, private answers and details, contacts & High\\
Identity & Verification, reset, enumeration and rate limits & High\\
Compatibility & SQLite locally, PostgreSQL 17 in CI, migrations on populated data & Medium--High\\
Usability & Multi-viewport sweep, keyboard flow, reduced motion & Medium--High\\
Model quality & Stability invariants and evaluation parity; real benchmark pending & Preliminary\\
\bottomrule
\end{tabularx}
\end{table}

% ===========================================================================
\chapter{Results and Discussion}
\textbf{In this chapter:}
\begin{itemize}[beginpenalty=10000]
  \item what was delivered against the synopsis;
  \item a worked matching example and the results on the demonstration data;
  \item pilot evaluation, performance observations and limitations.
\end{itemize}

\section{Application Output}
\begin{itemize}
  \item A moderated public wall and a personal home page.
  \item Explainable match views with evidence coverage, private claim timelines and the code handover.
  \item Notifications, staff review queues and administrator analytics.
  \item All screenshots come from the running server and a populated database, not wireframes.
\end{itemize}

\section{Features Successfully Implemented}
\begin{itemize}
  \item \textbf{Synopsis scope:} roles, report lifecycle, AJAX discovery, hybrid matching, explanations, claims, notifications, analytics, exports and responsive UI.
  \item \textbf{Added since the synopsis:} a stable versioned text model with titles, calibrated image similarity, colour sets, graded chronology, evidence coverage, retired match history, ownership reservation, disputes and reconsideration, the one-time-code handover, three-way themes, disc-overlap visualisation, the animated landing page and phone navigation.
  \item \textbf{Security and identity:} campus email verification, password recovery, rate limits, independent review, sanitised uploads, blurred sensitive photos, deny-by-default media, append-only audit records, and database constraints for one reserved claim, one open dispute and one active code.
\end{itemize}

\section{Worked Matching Example}\label{sec:worked}
\begin{itemize}
  \item Pair: ``JIIT student ID card'' (lost) and ``Student ID card found on a desk'' (found), shown in Figure~\ref{fig:explain}.
  \item The photographs show the same card, one clean and one worn; ResNet-50 gives a raw cosine of about 0.857, which calibration maps to $(0.857-0.275)/(1-0.275) \approx 0.80$.
  \item Title and description now share concepts (student ID card, JIIT, blue and white, B.Tech CSE), so the text score rose from 0.819 under TF--IDF to 0.936. Every signal had evidence on both reports, so coverage is 1.0 and no weight is redistributed.
\end{itemize}
\begin{table}[H]
\centering\small
\caption[Score breakdown for the ID-card pair]{Score breakdown for the ID-card pair}
\label{tab:worked}
\begin{tabularx}{\textwidth}{lccc X}
\toprule
Signal & $s_i$ & $w_i$ & $w_i s_i$ & Evidence\\
\midrule
Title + description (lexical) & 0.936 & 0.30 & 0.281 & Shared concepts: student ID card, JIIT, blue, white, B.Tech CSE\\
Photograph (ResNet-50) & 0.803 & 0.30 & 0.241 & Calibrated from raw cosine 0.857\\
Campus zone & 1.000 & 0.15 & 0.150 & Both Academic Block A\\
Category & 1.000 & 0.10 & 0.100 & Both ID \& Cards\\
Date and time & 1.000 & 0.10 & 0.100 & Same day (4 Nov); chronology plausible\\
Brand / colour & 1.000 & 0.05 & 0.050 & Both blue\\
\midrule
\textbf{Total} & & 1.00 & \textbf{0.922} & \textbf{Strong} ($>0.85$), coverage 1.0\\
\bottomrule
\end{tabularx}
\end{table}

\section{Matching on the Demonstration Data}
\begin{itemize}
  \item The chart breaks every live match into the points each signal contributed.
  \item Text and image vary most between pairs; zone, category, date and colour are near their maximum for all six.
\end{itemize}
\begin{figure}[!htbp]
\centering
\begin{tikzpicture}[x=13.2cm,y=0.72cm]
\def\pairs{{"JIIT ID card","Water bottle","Keys, red tag","Leather wallet","iPhone 13","Navy backpack"}}
\foreach \i/\t/\im/\l/\c/\d/\a in {0/0.2807/0.2410/0.15/0.1/0.1/0.05, 1/0.2908/0.1795/0.15/0.1/0.0933/0.05, 2/0.2237/0.3/0.075/0.1/0.0933/0.05, 3/0.2672/0.1609/0.15/0.1/0.0933/0.05, 4/0.2298/0.1192/0.15/0.1/0.1/0.05, 5/0.2513/0.0622/0.15/0.1/0.0933/0.05}{
  \pgfmathsetmacro{\xa}{\t}\pgfmathsetmacro{\xb}{\xa+\im}\pgfmathsetmacro{\xc}{\xb+\l}\pgfmathsetmacro{\xd}{\xc+\c}\pgfmathsetmacro{\xe}{\xd+\d}\pgfmathsetmacro{\xf}{\xe+\a}
  \fill[llaccent] (0,-\i-0.32) rectangle (\xa,-\i+0.32);
  \fill[llaccentink!55] (\xa,-\i-0.32) rectangle (\xb,-\i+0.32);
  \fill[teal!70] (\xb,-\i-0.32) rectangle (\xc,-\i+0.32);
  \fill[teal!35] (\xc,-\i-0.32) rectangle (\xd,-\i+0.32);
  \fill[llfound] (\xd,-\i-0.32) rectangle (\xe,-\i+0.32);
  \fill[llfound!45] (\xe,-\i-0.32) rectangle (\xf,-\i+0.32);
  \pgfmathsetmacro{\name}{\pairs[\i]}
  \node[font=\scriptsize,anchor=east] at (-0.01,-\i) {\name};
  \node[font=\scriptsize\bfseries,anchor=west] at (1.02,-\i) {\pgfmathprintnumber[fixed,precision=3]{\xf}};}
\node[font=\scriptsize\bfseries,anchor=west] at (1.02,0.75) {Score};
\foreach \x/\t in {0.70/moderate,0.85/strong} {\draw[ink,dashed,thick] (\x,0.55) -- (\x,-5.55); \node[font=\tiny,above] at (\x,0.55) {\t\ $\geq$ \x};}
\draw[ink!40] (0,0.55) -- (0,-5.55);
\foreach \x in {0,0.2,0.4,0.6,0.8,1.0} \node[font=\tiny,below] at (\x,-5.55) {\x};
\foreach \c/\t/\k in {llaccent/Text/0,llaccentink!55/Image/1,teal!70/Zone/2,teal!35/Category/3,llfound/Date/4,llfound!45/Attributes/5}
  {\fill[\c] (0.02+0.16*\k,-6.5) rectangle ++(0.025,0.36); \node[font=\scriptsize,anchor=west] at (0.05+0.16*\k,-6.32) {\t};}
\end{tikzpicture}
\caption[Score contributions of the live matches]{Contribution of each signal to the six demonstration matches (engine \code{synopsis-v3}); dashed lines are the heuristic tier thresholds}
\end{figure}

\begin{table}[H]
\centering\small
\caption[Matches on the demonstration database]{Matches on the demonstration database: engine synopsis-v2 versus synopsis-v3}
\label{tab:live}
\begin{tabularx}{\textwidth}{X X ccc}
\toprule
Lost report & Found report & Text (v2 $\rightarrow$ v3) & Image & Score v2 $\rightarrow$ v3\\
\midrule
JIIT student ID card & Student ID card found on a desk & 0.82 $\rightarrow$ 0.94 & 0.80 & 0.887 $\rightarrow$ \textbf{0.922} strong\\
Silver steel water bottle & Metal water bottle, steel colour & 0.87 $\rightarrow$ 0.97 & 0.60 & 0.835 $\rightarrow$ \textbf{0.864} strong\\
Bunch of keys with red tag & Keys with a red keychain & 0.62 $\rightarrow$ 0.75 & 1.00 & 0.805 $\rightarrow$ \textbf{0.842} moderate\\
Black leather wallet & Black wallet made of leather with silver logo & 0.84 $\rightarrow$ 0.89 & 0.54 & 0.807 $\rightarrow$ \textbf{0.821} moderate\\
iPhone 13 with blue case & Blue Apple phone found on table & 0.52 $\rightarrow$ 0.77 & 0.40 & 0.675 $\rightarrow$ \textbf{0.749} moderate\\
Navy blue backpack & Dark blue bag left in classroom & 0.66 $\rightarrow$ 0.84 & 0.21 & 0.654 $\rightarrow$ \textbf{0.707} moderate\\
\bottomrule
\end{tabularx}
\end{table}

\begin{itemize}
  \item \textbf{All six matches are true pairs}, and no wrong pairing reaches the 0.40 storage threshold under either engine.
  \item \textbf{Only the text component changed} between v2 and v3; image, zone, category, date and colour scores are identical, and every pair has full evidence coverage, no tier cap and no chronology warning.
  \item \textbf{Titles and synonyms matter most where wording differs:} ``iPhone 13'' versus ``Blue Apple phone'' and ``backpack'' versus ``bag'' gained the most and moved from weak to moderate; they stay visible but are clearly below the strong pairs.
  \item \textbf{The keys pair} reuses one demo photograph for both reports, so its perfect image score is an artefact of the demo data.
  \item \textbf{Stability:} rebuilding twice changed nothing, and adding unrelated reports leaves every score unchanged, which the v2 text model could not guarantee.
  \item \textbf{Thresholds unchanged:} the synopsis tiers were not retuned and are labelled heuristic, because the available evaluation data cannot justify new ones.
\end{itemize}

\section{Pilot Evaluation}
\begin{itemize}
  \item 180 pairs from 60 Wikimedia Commons images in six categories, with a source-disjoint validation and test split.
  \item Results below were regenerated for \code{synopsis-v3} with the same scoring function as production; thresholds were chosen on the validation split only. The benchmark is optimistic: positives reuse the same photo with brightness or compression changes, and descriptions come from source titles.
  \item Each results file records the dataset hash and provenance, engine and model versions, weights, thresholds and limitations, and is stamped \emph{engineering check only}.
  \item They show the pipeline works end to end; they are \textbf{not} evidence of real-world accuracy, and the set intentionally fails the project's research-readiness gate.
\end{itemize}
\begin{table}[H]
\centering\small
\caption[Pilot test-split results]{Pilot test-split results (engine \code{synopsis-v3}; engineering sanity check only)}
\label{tab:pilot}
\footnotesize
\begin{tabular}{llcccccc}
\toprule
Text & Image & Threshold & Precision & Recall & F1 & MRR & F1, no image\\
\midrule
Lexical & ResNet-50 & 0.75 & 0.833 & 1.000 & 0.909 & 1.000 & 0.759\\
Lexical & CLIP & 0.80 & 0.909 & 1.000 & 0.952 & 1.000 & 0.779\\
Lexical & none & 0.95 & 0.909 & 1.000 & 0.952 & 1.000 & 0.952\\
\bottomrule
\end{tabular}
\end{table}

\section{Performance and Observations}
\begin{itemize}
  \item Feature caching avoids recomputing vectors for unchanged reports; identical reruns write nothing, so query counts stay flat as the pool grows.
  \item Pagination bounds every list; the lexical text model is a small JSON release and costs almost nothing at campus scale.
  \item ResNet-50 on a CPU adds noticeable latency at approval, because matching runs synchronously in the request.
  \item The interface shows the score components immediately once the request completes.
\end{itemize}
\begin{table}[H]
\centering\small
\caption{Observed strengths and evidence}
\begin{tabularx}{\textwidth}{P{32mm}X}
\toprule
Strength & Evidence\\
\midrule
Traceability & Components, evidence status, coverage, chronology, models and versions are stored; history is retired, not deleted\\
Honest scoring & Stable text scores, calibrated images, partial colours and graded chronology; scores labelled as similarity, not probability\\
Human control & Similarity never approves ownership; an independent moderator decides with a checklist\\
Privacy & Answers, private details, evidence, contacts and sensitive photos scoped to those who need them\\
Case integrity & Reservation survives disputes; only the owner's code (or a documented exception) returns an item\\
Design coherence & Tokens, overlapping match discs and status dots form one visual language\\
\bottomrule
\end{tabularx}
\end{table}

\section{Limitations}
\begin{itemize}
  \item Real-world accuracy is unmeasured; the pilot set reuses source photographs, and the chance levels were measured on that same set.
  \item The text model's class weights, the coverage caps and the chronology thresholds are documented heuristics, not calibrated statistics.
  \item One photograph per report; matching runs synchronously; automatic sensitive-photo detection is only a keyword hint for moderators.
  \item Production email needs an SMTP account; notifications other than account email are in-app only; no native mobile client.
  \item Only a free demonstration instance has been deployed; the production deployment is documented but not launched.
\end{itemize}

% ===========================================================================
\chapter{Conclusion and Future Scope}
\textbf{In this chapter:}
\begin{itemize}[beginpenalty=10000]
  \item the main conclusions of the project;
  \item a prioritised roadmap of future enhancements;
  \item smaller improvements for operation in production.
\end{itemize}

\section{Conclusion}
\begin{itemize}
  \item LostLink delivers a complete, responsible campus recovery workflow: structured reports, explainable multimodal ranking and human-controlled verification.
  \item Its central design decision, separating a \emph{probable match} from an \emph{approved owner}, allows useful automation without giving up privacy or accountability.
  \item The remediation made the system more honest and harder to misuse: pair scores are now stable and labelled as similarity, sparse evidence is visible as coverage, and a dispute can no longer release an item to someone else.
  \item Claim approved is no longer item returned: only the owner's one-time code, shared after they hold the item, closes a case.
  \item Two logic audits closed eighteen gaps where stored state could contradict reality, each now guarded by a test.
  \item The codebase is modular, tested (364 tests plus browser, mobile and first-run checks) and deployment-ready, and it states its evidence limitations openly.
\end{itemize}

\section{Future Enhancements}
\begin{itemize}
  \item Priorities reflect what limits the project most today: evidence of accuracy, then a live deployment.
\end{itemize}
\begin{table}[!tbp]
\centering\small
\caption{Future work roadmap}
\begin{tabularx}{\textwidth}{P{6mm}X P{18mm}P{30mm}}
\toprule
\# & Enhancement & Priority & Benefit\\
\midrule
1 & Manually reviewed dataset of independently photographed campus items; re-measure chance levels; report precision, recall, F1, MRR and recall@K before retuning thresholds & High & Real accuracy evidence\\
2 & Launch the PostgreSQL deployment with SMTP and a shared cache, and run \code{rebuild\_matches} with every engine upgrade & High & Live service\\
3 & Background queue for matching and feature extraction, with progress and retries & High & Faster approvals\\
4 & Multiple photographs per report, with perceptual-hash duplicate detection & Medium & Better image evidence\\
5 & Optional email or SMS alerts for matches and claims, with preferences & Medium & Faster response\\
6 & Object storage, antivirus scanning and automated redaction suggestions for sensitive photos & Medium & Safer media\\
7 & Run the documented AJAX-versus-manual usability study with consenting participants & Medium & Usability evidence\\
8 & Progressive web app or native client once the workflow stabilises & Low & Offline, app-like use\\
\bottomrule
\end{tabularx}
\end{table}

\section{Possible Improvements}
\begin{itemize}
  \item Query profiling and asynchronous exports.
  \item A full accessibility audit.
  \item An administrator view comparing engine versions side by side.
  \item Monitoring that pairs score-drift snapshots with periodically reviewed ground truth, since drift alone cannot measure correctness.
  \item Production operations: backups, structured logs, error tracking, HTTPS and an incident process.
\end{itemize}

% ===========================================================================
\chapter{References}
\begingroup
\fontsize{9.5}{11.4}\selectfont
\begin{thebibliography}{99}
\setlength{\itemsep}{0.5pt}
\setlength{\parskip}{0pt}

\bibitem{python}
Python Software Foundation, ``Python 3 documentation,'' Python.org. Accessed: Nov. 11, 2026. [Online]. Available: \url{https://docs.python.org/3/}

\bibitem{django}
Django Software Foundation, ``Django documentation,'' DjangoProject.com. Accessed: Nov. 11, 2026. [Online]. Available: \url{https://docs.djangoproject.com/}

\bibitem{pillow}
A. Clark and contributors, ``Pillow (PIL Fork) documentation,'' Read the Docs. Accessed: Nov. 11, 2026. [Online]. Available: \url{https://pillow.readthedocs.io/}

\bibitem{sklearn}
scikit-learn developers, ``scikit-learn user guide: Feature extraction and pairwise metrics,'' scikit-learn.org. Accessed: Nov. 11, 2026. [Online]. Available: \url{https://scikit-learn.org/stable/}

\bibitem{sklearnpaper}
F. Pedregosa \emph{et al.}, ``Scikit-learn: Machine learning in Python,'' \emph{J. Mach. Learn. Res.}, vol. 12, pp. 2825--2830, 2011.

\bibitem{reportlab}
ReportLab Inc., ``ReportLab PDF library user guide,'' ReportLab.com. Accessed: Nov. 11, 2026. [Online]. Available: \url{https://docs.reportlab.com/}

\bibitem{pytorch}
PyTorch Foundation, ``PyTorch documentation,'' PyTorch.org. Accessed: Nov. 11, 2026. [Online]. Available: \url{https://pytorch.org/docs/stable/}

\bibitem{pytorchpaper}
A. Paszke \emph{et al.}, ``PyTorch: An imperative style, high-performance deep learning library,'' in \emph{Proc. Adv. Neural Inf. Process. Syst. (NeurIPS)}, vol. 32, 2019, pp. 8024--8035.

\bibitem{postgres}
PostgreSQL Global Development Group, ``PostgreSQL 17 documentation,'' PostgreSQL.org. Accessed: Nov. 11, 2026. [Online]. Available: \url{https://www.postgresql.org/docs/17/}

\bibitem{mdn}
Mozilla Developer Network, ``MDN Web Docs: HTML, CSS and JavaScript,'' developer.mozilla.org. Accessed: Nov. 11, 2026. [Online]. Available: \url{https://developer.mozilla.org/}

\bibitem{whitenoise}
D. Evans, ``WhiteNoise documentation,'' Read the Docs. Accessed: Nov. 11, 2026. [Online]. Available: \url{https://whitenoise.readthedocs.io/}

\bibitem{tfidf}
G. Salton and C. Buckley, ``Term-weighting approaches in automatic text retrieval,'' \emph{Inf. Process. Manage.}, vol. 24, no. 5, pp. 513--523, 1988, doi: 10.1016/0306-4573(88)90021-0.

\bibitem{resnet}
K. He, X. Zhang, S. Ren, and J. Sun, ``Deep residual learning for image recognition,'' in \emph{Proc. IEEE Conf. Comput. Vis. Pattern Recognit. (CVPR)}, Las Vegas, NV, USA, 2016, pp. 770--778, doi: 10.1109/CVPR.2016.90.

\bibitem{clip}
A. Radford \emph{et al.}, ``Learning transferable visual models from natural language supervision,'' in \emph{Proc. 38th Int. Conf. Mach. Learn. (ICML)}, 2021, pp. 8748--8763.

\bibitem{lostlinkrepo}
S. Grover, P. Goel, A. Sharma, and A. Kulashri, ``LostLink: An AI-enabled lost-and-found matching platform,'' GitHub repository, 2026. Accessed: Nov. 11, 2026. [Online]. Available: \url{https://github.com/SaiGrover/LostLink}

\bibitem{sbert}
N. Reimers and I. Gurevych, ``Sentence-BERT: Sentence embeddings using Siamese BERT-networks,'' in \emph{Proc. Conf. Empirical Methods Natural Lang. Process. (EMNLP-IJCNLP)}, Hong Kong, China, 2019, pp. 3982--3992, doi: 10.18653/v1/D19-1410.

\bibitem{playwright}
Microsoft, ``Playwright for Python documentation,'' Playwright.dev. Accessed: Nov. 11, 2026. [Online]. Available: \url{https://playwright.dev/python/}

\end{thebibliography}
\endgroup

% ===========================================================================
\chapter*{Appendix A: Setup and Commands}
\addcontentsline{toc}{chapter}{Appendix A: Setup and Commands}
\markboth{Appendix A}{Appendix A}
\textbf{Run LostLink locally (demo data)}
\begin{lstlisting}[style=lostlink,language=bash]
python -m venv .venv
.venv\Scripts\activate            # macOS/Linux: source .venv/bin/activate
pip install -r requirements-ml.txt
python manage.py migrate
python manage.py seed_demo --fresh
python manage.py upgrade_demo_photos  # credited demo photos; skip offline
python manage.py rebuild_matches      # downloads pinned ResNet-50 weights once
python manage.py runserver            # open http://127.0.0.1:8000/
\end{lstlisting}
\textbf{Verify the build}
\begin{lstlisting}[style=lostlink,language=bash]
python manage.py check
python manage.py makemigrations --check --dry-run
python manage.py test                 # 364 tests; lexical text + histogram, no downloads
python scripts/smoke_first_run.py     # empty database to code-confirmed handover
python scripts/browser_check.py --base http://127.0.0.1:8000
python scripts/e2e_workflow.py --base http://127.0.0.1:8000
python scripts/mobile_nav_check.py --base http://127.0.0.1:8000
\end{lstlisting}
\textbf{Notes}
\begin{itemize}
  \item Demo accounts are created by \code{seed\_demo} and listed on the sign-in page only while \code{DEBUG} is on; never seed a production database.
  \item After changing the engine version or image model, run \code{rebuild\_matches} once; reports and claims are not touched.
  \item Production uses PostgreSQL via \code{DATABASE\_URL}, \code{DJANGO\_DEBUG=false}, a private secret key, exact hosts, SMTP settings (\code{DJANGO\_EMAIL\_*}) and \code{seed\_taxonomy --if-empty} instead of demo data (see \code{docs/DEPLOY\_RENDER.md}).
  \item Scheduled maintenance: \code{purge\_claim\_evidence} daily; \code{retry\_features} after model outages.
\end{itemize}

\setcounter{table}{0}\renewcommand{\thetable}{B.\arabic{table}}
\chapter*{Appendix B: Repository Layout}
\addcontentsline{toc}{chapter}{Appendix B: Repository Layout}
\markboth{Appendix B}{Appendix B}
\begin{itemize}
  \item Each Django app sits in its own top-level folder with its models, views, URLs, forms and tests.
  \item Templates, static files and documentation are shared across the apps.
\end{itemize}
\begin{table}[H]
\centering\footnotesize\renewcommand{\arraystretch}{1.0}
\caption{Top-level folders of the LostLink repository}
\begin{tabularx}{\textwidth}{P{34mm}X}
\toprule
Folder & Contents\\
\midrule
\code{config/} & Django settings, root URL configuration and WSGI entry point\\
\code{accounts/} & User model, roles, verification, password reset, rate limiter, decorators\\
\code{items/} & Reports, lifecycle, search, media security, taxonomy and demo commands\\
\code{matching/} & Text model release, features, calibration, chronology, scoring, evaluation\\
\code{claims/} & Claim services, disputes, revisions, evidence, handover codes\\
\code{notifications/} & In-app notifications and the unread-count endpoint\\
\code{dashboard/} & Review queue, decision log, users, taxonomy, analytics and exports\\
\code{templates/} & Shared page shell, partials and per-app page templates\\
\code{static/} & Stylesheets (with the design tokens) and JavaScript\\
\code{scripts/} & First-run smoke test, browser sweep, mobile check, end-to-end workflow, PostgreSQL helper\\
\code{evaluation/} & Pilot dataset, evaluation results and study templates\\
\code{docs/} & Claims and handover, deployment, evaluation, frontend and PostgreSQL notes\\
\code{report/} & This report's LaTeX source, figures and PDFs\\
\bottomrule
\end{tabularx}
\end{table}

\end{document}

